% !TEX root = immunology.tex

% \section{Immune System Players}

%   \begin{core}[label=core:hsc]{Hematopoietic Stem Cells (HSCs) \& Lineage}

%     \textbf{Hematopoietic stem cells (HSCs)} are pluripotent stem cells capable of giving rise to any type of blood cell. Each HSC can either divide to produce additional copies of itself---a process termed self-renewal---or begin to differentiate through a series of lineage-restricting choices.

%     Developmentally, HSCs first arise in the yolk sac membrane of the early embryo, subsequently migrate to the liver and spleen, and ultimately settle predominantly in the \textbf{bone marrow} prior to birth. Remarkably, as few as approximately 100 HSCs are sufficient to fully regenerate the entire hematopoietic system.

%     \autoref{fig:hsc_lineage} illustrates the hierarchical differentiation of HSCs into the myeloid and lymphoid lineages, highlighting the key transcription factors governing lineage commitment at each branch point.

%     \pgfig{
%       % Styles
%       \tikzset{
%         myeloid/.style={rectangle, draw, fill={rgb,255:red,164;green,194;blue,244}, minimum width=3cm, minimum height=1cm, align=center, font=\small},
%         lymphoid/.style={rectangle, draw, fill={rgb,255:red,249;green,203;blue,156}, minimum width=2cm, minimum height=1cm, align=center, font=\small},
%         leafnode/.style={rectangle, rounded corners=8pt, draw, fill={rgb,255:red,207;green,226;blue,243}, minimum width=2.4cm, minimum height=0.9cm, align=center, font=\small},
%         edgelabel/.style={font=\small\itshape, align=left}
%       }
%       % Level 0
%       \node[myeloid] (HSC) at (5,0) {Haematopoietic\\Stem Cell};
%       % Level 1
%       \node[myeloid]  (CMP) at (1,-2)   {Common Myeloid\\Progenitor (CMP)};
%       \node[lymphoid] (CLP) at (9,-2)   {Common Lymphoid\\Progenitor (CLP)};
%       % Level 2
%       \node[myeloid]  (MEP) at (-1.5,-4) {Megakaryocyte-Erythroid\\Progenitor (MEP)};
%       \node[myeloid]  (GMP) at (2.5,-4)  {Granulocyte-Macrophage\\Progenitor (GMP);\\Myeloblast};
%       \node[lymphoid] (NKP) at (6,-4)  {NK Cell\\Progenitor};
%       \node[lymphoid] (BCP) at (9,-4)  {B Cell\\Progenitor};
%       \node[lymphoid] (TCP) at (12,-4)   {T Cell\\Progenitor};
%       % Level 3 - GMP lineage
%       \node[leafnode] (BASO) at (-3,-7.8) {Basophil};
%       \node[leafnode] (EOSI) at (-0,-7.8) {Eosinophil};
%       \node[leafnode] (CDC)  at (3,-7.8) {cDC};
%       \node[leafnode] (NEUT) at (6,-7.8) {Neutrophil};
%       \node[leafnode] (MONO) at (1,-9.4) {Monocyte};
%       \node[leafnode] (MAC)  at (-0.5,-11.0) {Macrophage};
%       \node[leafnode] (MODC) at (2.5,-11.0) {moDCs};
%       % Arrows: GMP
%       \draw[-{latex}, thick] (GMP) -- (BASO);
%       \draw[-{latex}, thick] (GMP) -- (EOSI);
%       \draw[-{latex}, thick] (GMP) -- (CDC);
%       \draw[-{latex}, thick] (GMP) -- (NEUT);
%       \draw[-{latex}, thick] (GMP) -- (MONO);
%       \draw[-{latex}, thick] (MONO) -- (MAC);
%       \draw[-{latex}, thick] (MONO) -- (MODC);
%       % Level 3 - MEP lineage
%       \node[leafnode] (RBC)  at (-3,-6.2) {RBCs\\(Erythrocytes)};
%       \node[leafnode] (MEGA) at (-0.5,-6.2) {Megakaryocytes;\\Platelets\\(Thrombocytes)};
%       % Level 3 - lymphoid leaves
%       \node[leafnode] (NKC) at (6.3,-6.2)  {NK Cells};
%       \node[leafnode] (BC)  at (9,-6.2)  {B Cells};
%       \node[leafnode] (TCC) at (13,-6.2)   {TC Cells};
%       % T helper branch
%       \node[leafnode] (THC)   at (10.6,-7.8) {TH Cells};
%       \node[leafnode] (TH2)   at (9.5,-10.8) {TH2 Cells};
%       \node[leafnode] (TH17)  at (12.1,-10.8) {TH17 Cells};
%       \draw[-{latex}, thick] (THC) -- (TH2);
%       \draw[-{latex}, thick] (THC) -- (TH17);
%       \node[leafnode] (TH1)   at (8.6,-9.6) {TH1 Cells};
%       \node[leafnode] (THREG) at (13.0,-9.6) {THreg Cells};
%       \draw[-{latex}, thick] (THC) -- (TH1);
%       \draw[-{latex}, thick] (THC) -- (THREG);
%       % Arrows: HSC
%       \draw[-{latex}, thick] (HSC) -- (CMP);
%       \draw[-{latex}, thick] (HSC) -- (CLP);
%       \draw[-{latex}, thick] (CMP) -- (MEP);
%       \draw[-{latex}, thick] (CMP) -- (GMP);
%       \draw[-{latex}, thick] (CLP) -- (NKP);
%       \draw[-{latex}, thick] (CLP) -- (BCP);
%       \draw[-{latex}, thick] (CLP) -- (TCP);
%       % Arrows: MEP
%       \draw[-{latex}, thick] (MEP) -- (RBC);
%       \draw[-{latex}, thick] (MEP) -- (MEGA);
%       % Arrows: lymphoid leaves
%       \draw[-{latex}, thick] (NKP) -- (NKC);
%       \draw[-{latex}, thick] (BCP) -- (BC);
%       \draw[-{latex}, thick] (TCP) -- (TCC);
%       \draw[-{latex}, thick] (TCP) -- (THC);
%       % Standalone edge labels (fixed coordinates, independent of node positions)
%       \node[edgelabel] at (5,0.9)   {\underline{\textbf{Bmi-1}} to stay in renewal cycle (Lin-)};
%       \node[edgelabel] at (9,-0.9) {\underline{\textbf{Ikaros}} transcription\\factor to become CLP};
%       \node[edgelabel] at (5,-2) {\underline{\textbf{GATA-2}} to\\differentiate (Lin+)};
%       \node[edgelabel] at (12.5,-2.9) {\underline{\textbf{Notch family}} for T cell\\lineage pathway};
%     }{Hierarchy of haematopoietic stem cell differentiation into myeloid (\textit{left}) and lymphoid (\textit{right}) lineages}{fig:hsc_lineage}{0.85}

%     As shown in \autoref{fig:hsc_lineage}, HSCs destined for differentiation are termed \textbf{lineage positive (\textit{lin+})}, whereas those not destined for differentiation are termed \textbf{lineage negative (\textit{lin-})}. This binary state is governed by the presence of one of two transcription factors.

%     \begin{itemize}

%       \item \textbf{Bmi-1}: A transcription factor whose presence maintains HSCs in an \textbf{undifferentiated state}, sustaining ongoing self-renewal and division.

%       \item \textbf{GATA-2}: A transcription factor that \textbf{triggers HSC differentiation}, initiating the developmental path toward a specialized cell type.

%     \end{itemize}

%     Given a \textit{lin+} state, the first branch point determines whether the cell commits to the myeloid or lymphoid lineage.

%     \begin{itemize}

%       \item Commitment to a \textbf{common myeloid-erythroid progenitor} enables development into a wide array of cell types, including various white blood cells, red blood cells, and platelet-producing megakaryocytes.

%       \item Commitment to a \textbf{common lymphoid progenitor} is marked by expression of the transcription factor \textbf{Ikaros}, giving rise to adaptive T and B cells or innate NK and dendritic cells. Further along the differentiation tree, the \textbf{Notch family} of transcription factors governs the divergence of T and B cell development.

%     \end{itemize}

%   \end{core}

%   \begin{core}{Lymphoid Lineage}

%     Cells of the \textbf{lymphoid lineage} constitute 20--40\% of circulating leukocytes in the blood and approximately 99\% of the cells of the lymph. These cells are broadly categorized as either \textbf{adaptive}---characterized by gene rearrangement during the immune response---or \textbf{innate}---characterized by the absence of such gene rearrangement.

%     \autoref{tab:lymphocytes} presents a structured overview of the cell types comprising the lymphoid lineage depicted in \autoref{fig:hsc_lineage}.

%   \end{core}

%   \begin{sidewaystable}

%     \centering
%     \caption{Overview of adaptive and innate lymphocyte populations}
%     \label{tab:lymphocytes}
%     \begin{tabular}{
%       >{\raggedright\arraybackslash}p{1.2cm}
%       >{\raggedright\arraybackslash}p{2.2cm}
%       >{\raggedright\arraybackslash}p{3.0cm}
%       >{\raggedright\arraybackslash}p{3.8cm}
%       >{\raggedright\arraybackslash}p{2.5cm}
%       >{\raggedright\arraybackslash}p{2.5cm}
%       >{\raggedright\arraybackslash}p{3.3cm}
%     }
%     \hline
%     \textbf{Type} & \textbf{Name} & \textbf{Function} & \textbf{Surface Expression} & \textbf{Responds to} & \textbf{Releases} & \textbf{Notes} \\
%     \hline

%     \multirow{5}{*}{\textbf{ADAPTIVE}}
%       % =====================================================================================
%       & \textit{T$_\text{H}$ (Helper) Lymphocyte} & \textbf{Activates} inflammation, macrophages, and B cells
%       & [CD4]---{\footnotesize stabilizes MHC~II--TCR antigen recognition}

%       [TCR]
%       & MHC~II-presented antigen & Cytokines
%       & \footnotesize
%         Subtypes include:

%         \textbf{T$_\text{H}$1}---direct response

%         \textbf{T$_\text{H}$2}---nutrient restraint response

%         \textbf{T$_\text{H}$17}---fungi/parasite response

%         \textbf{Memory cells}---future reserve \\[4pt]

%       % =====================================================================================
%       & \textit{T$_\text{R}$ (Regulatory) Lymphocyte} & \textbf{Suppresses} immune response (\textit{generally and of specific lymphocytes})
%       & [CD4]

%       [CD25]---{\footnotesize stabilizes T$_\text{reg}$ binding to IL-2, essential for survival and function}

%       & \textit{Unknown}
%       & Cytokines;

%       IL-10;

%       TGF-β
%       & \footnotesize Prevents allergy and autoimmune disease; subtype of the TH lineage \\[4pt]

%       % =====================================================================================
%       & \textit{T$_\text{C}$ (Cytotoxic) Lymphocyte} & \textbf{Destroys} altered-self cells
%       & [CD8]---{\footnotesize stabilizes MHC~I--TCR antigen recognition}
%       & MHC~I-presented antigen
%       & Perforin;

%       Granzymes
%       & \\[4pt]

%       % =====================================================================================
%       & \textit{B Cell} & \textbf{Recognizes} antigen; can present antigen to T cells
%       & [Membrane-bound antibodies/BCRs]

%       [MHC~II]---{\footnotesize constitutive; presents BCR-captured antigen fragments}

%       [TLR]

%       [Fc receptors]---{\footnotesize down-regulate antibody production}

%       [B7]
%       & Circulating antigen;

%       T$_\text{H}$ cytokine signals & [B7]
%       & \footnotesize Develops into a \textbf{plasma cell} or \textbf{memory cell} (\textit{reserve for future infection}) \\[4pt]

%       % =====================================================================================
%       & \textit{Plasma Cell} & \textbf{Secretes} antibodies
%       & No surface antibodies & T$_\text{H}$ cytokine signals & Antibodies
%       & \footnotesize Mature B cell;

%       Lives 1--2 weeks \\
%     \hline

%     \multirow{1}{*}{\textbf{INNATE}}
%       % =====================================================================================
%       & \textit{NK (Natural Killer) Cell} & \textbf{Destroys} altered-self cells
%       & [CD16]---{\footnotesize Fc receptor recognizing antibody-coated altered cells, triggering attack}

%       [MHC~I receptors]---{\footnotesize induce killing when MHC~I is absent}
%       & Antibodies;

%       downregulated MHC~I 
%       & Perforin;

%       Granzymes
%       & \footnotesize Executes \textbf{antibody-dependent cell-mediated cytotoxicity (ADCC)}, triggering apoptosis via antibody signals on antigen-coated cell surfaces \\

%     \hline
%     \end{tabular}

%   \end{sidewaystable}

%   \begin{core}{Myeloid Lineage}

%     Cells of the \textbf{myeloid lineage} can be broadly categorized into \textbf{granulocytic cells}---characterized by specific granules that compartmentalize potentially dangerous molecules---and \textbf{antigen-presenting cells (APCs)}, which present antigen to T cells. Additional myeloid lineage cell types include \textbf{red blood cells (\textit{erythrocytes})} and \textbf{platelets (\textit{thrombocytes})}, which serve more secondary roles within general immune function.

%     \autoref{tab:myeloid} presents a structured overview of the cell types comprising the myeloid lineage depicted in \autoref{fig:hsc_lineage}.

%   \end{core}

%   \begin{sidewaystable}

%     \centering
%     \caption{Overview of innate granulocytic (\textit{1st page}), innate antigen presenting (\textit{2nd page}), and other myeloid cell populations (\textit{2nd page})}
%     \label{tab:myeloid}
%     \begin{tabular}{
%       >{\raggedright\arraybackslash}p{1.2cm}
%       >{\raggedright\arraybackslash}p{1.5cm}
%       >{\raggedright\arraybackslash}p{3.0cm}
%       >{\raggedright\arraybackslash}p{4.7cm}
%       >{\raggedright\arraybackslash}p{2.2cm}
%       >{\raggedright\arraybackslash}p{2.5cm}
%       >{\raggedright\arraybackslash}p{3.3cm}
%     }
%     \hline
%     \textbf{Type} & \textbf{Name} & \textbf{Function} & \textbf{Surface Expression} & \textbf{Responds to} & \textbf{Releases} & \textbf{Notes} \\
%     \hline

%     \multirow{5}{1.2cm}{\textbf{INNATE (Granul- ocytic)}}
%       % =====================================================================================
%       & \textit{Neutrophil}
%       & \textbf{Phagocytosis}, mediates the acute inflammatory response

%       Migrates into infected tissues
%       & [Fc Receptors]---{\footnotesize binds to the \textbf{fragment crystallizable} region of antibodies, facilitating the recognition of antibody-coated pathogens and subsequent immune activation}

%       [Complement Receptors] ---{\footnotesize enhances the phagocytosis of pathogens opsonized by complement proteins}
      
%       [PRR]---{\footnotesize mediates pathogen pattern recognition; specifically \textbf{Toll-like receptors}}
%       & Chemokines
%       & Acidic \& basic granules

%       ROS/RNS via the NOX complex

%       [IL-1]

%       [IL-6]

%       [TNFα]
%       & \footnotesize
%       First responders to inflection

%       One-day lifespan

%       Multilobed nucleus \\[4pt]
%       % =====================================================================================
%       & \textit{Basophil}
%       & \textbf{Defense} against allergic responses and parasitic infections
%       & [Fc Receptors]---{\footnotesize for IgE antibodies}
%       & \textit{Chemokines}
%       & Histamine-rich basophilic granules
%       & \footnotesize Non-phagocytic

%       Lobed nucleus

%       Responds to worms\\[4pt]
%       % =====================================================================================
%       & \textit{Mast Cell}
%       & \textbf{Defense} against allergic responses and parasitic infections
%       & [Fc Receptors]---{\footnotesize for IgE antibodies}
%       & 
%       & Histamine-rich basophilic granules
%       & \footnotesize Non-phagocytic

%       Mononuclear (non-lobed) \\[4pt]
%       % =====================================================================================
%       & \textit{Eosinophil}
%       & \textbf{Defense} against allergic responses and parasitic infections
%       & [Fc Receptors]---{\footnotesize for IgE antibodies}
%       & \textit{Chemokines}
%       & Eosinophilic granules containing hydrolytic enzymes
%       & \footnotesize Phagocytic, though not its primary function

%       Bi-lobed nucleus
%       \\[4pt]
%       \hline
%       % =====================================================================================
%       &
%       &
%       &
%       &
%       &
%       & \raggedleft \footnotesize continued on next page \\

%     \end{tabular}

%   \end{sidewaystable}

%   \begin{sidewaystable}

%     \centering
%     \begin{tabular}{
%       >{\raggedright\arraybackslash}p{1.2cm}
%       >{\raggedright\arraybackslash}p{2.2cm}
%       >{\raggedright\arraybackslash}p{3.0cm}
%       >{\raggedright\arraybackslash}p{3.8cm}
%       >{\raggedright\arraybackslash}p{2.5cm}
%       >{\raggedright\arraybackslash}p{2.5cm}
%       >{\raggedright\arraybackslash}p{3.3cm}
%     }
%     \hline
%     \textbf{Type} & \textbf{Name} & \textbf{Function} & \textbf{Surface Expression} & \textbf{Responds to} & \textbf{Releases} & \textbf{Notes} \\
%     \hline

%     \multirow{3}{1.2cm}{\textbf{INNATE (Antigen Present- ing)}}
%       % =====================================================================================
%       & \textit{Sentinel Dendritic Cell (DC)} 
%       & \textbf{Processes and presents} antigens to T$_\text{H}$ cells

%       \textit{Minor roles:} \textbf{Phagocytosis}, \textbf{endocytosis}, and \textbf{pinocytosis} of microorganisms and cellular debris
%       & [MHC~II]---{\footnotesize constitutive}

%       [B7]---{\footnotesize facilitates T cell activation}

%       [PRR]---{\footnotesize mediates pathogen pattern recognition; specifically \textbf{Toll-like receptors}}
%       & 
%       & [B7]---{\footnotesize co-stimulatory molecule provided to T cells during MHC-TCR antigen presentation}
%       & \footnotesize
%       DCs can differentiate from CMPs, monocytes, and CLPs (\textit{each yielding distinct specialized subsets})

%       Reside in peripheral tissues, \textbf{migrating to secondary lymphoid organs} only upon detecting pathogenic signals \\[4pt]
%       % =====================================================================================
%       & \textit{Monocyte}
%       & \textbf{Precursors} to macrophages, DCs, and MDSCs
%       & [TLR]
%       & Interferon from T$_\text{H}$ cells
%       & 
%       & \footnotesize Circulates in the bloodstream for approximately 8 hours \\[4pt]

%       % =====================================================================================
%       & \textit{Macrophage}
%       & \textbf{Phagocytosis} of microorganisms and cellular debris

%       \textit{Minor role:} \textbf{Presents} antigens following activation
%       & [MHC~II]---{\footnotesize post-activation}

%       [PRR]---{\footnotesize mediates pathogen pattern recognition; specifically \textbf{Toll-like receptors}}
%       & 
%       & ROS/RNS via the NOX complex

%       [B7]

%       [IL-1],[IL-6]

%       [TNFα]
%       & \footnotesize Exhibits amoeboid motility

%       Pathogens opsonized by antibodies are more readily phagocytized\\[4pt]
%     \hline

%     \multirow{2}{*}{\textbf{OTHER}}
%       % =====================================================================================
%       & \textit{Red Blood Cells (RBC; Erythrocytes)} 
%       & \textbf{Distributes} oxygen and nutrients systemically
%       & [CD235a]---{\footnotesize contributes to RBC antigenic properties}
%       & 
%       & 
%       & \\[4pt]
%       % =====================================================================================
%       & \textit{Platelets (Thrombocytes)} 
%       & \textbf{Mediates} coagulation at sites of vascular injury
%       & [CD41]---{\footnotesize functions as a key marker essential for platelet aggregation}

%       [CD42b]---{\footnotesize facilitates platelet adhesion via a glycoprotein complex}
%       & 
%       & 
%       & \footnotesize Platelets are derived from megakaryocytes\\

%     \hline
%     \end{tabular}

%   \end{sidewaystable}

%   \begin{core}{Follicular Dendritic Cells}

%     \textbf{Follicular dendritic cells (FDCs)} are specialized stromal cells that, in contrast to the myeloid, antigen presenting \textbf{sentinel dendritic cells} shown in \autoref{tab:myeloid}, do not arise via \textbf{hematopoiesis} from HSCs. Instead, they originate during embryogenesis from a distinct \textbf{mesenchymal} lineage of \textbf{mesenchymal progenitor cells}.
    
%     Following development, FDCs migrate into lymphoid tissues, where they ultimately become embedded within the follicles of secondary lymphoid organs.
    
%     Once established, FDCs are distinguished by their unique capacity to \textbf{retain and present intact antigens} on their surfaces over extended periods. They capture antigens from the lymph and retain them as immune complexes, a function essential for \textbf{B cell activation} and \textbf{germinal center formation}.
    
%     The most characteristic surface markers of FDCs are as follows.

%     \begin{itemize}

%       \item \textbf{CD21/CD35}: FDCs express the complement receptors CD21 and CD35, which mediate their role in antigen capture.

%       \item \textbf{FDC-M1}: This marker is specific to FDCs, enabling their identification and distinction from other cellular populations within lymphoid tissue.

%     \end{itemize}

%   \end{core}

% \section{Immune System Organs}
  
%   \subsection{Primary Immune System Organs}

%     \begin{core}{Primary Immune System Organs}

%       The \textbf{primary immune system organs} are the sites at which cells divide, commit to a developmental fate, and, where applicable, undergo gene rearrangement.

%     \end{core}

%     \begin{core}{Bone Marrow}

%       \sbs{0.5}{
%         The \textbf{bone marrow} is the soft, highly vascularized tissue residing within the hollow interior cavities of bones. It is the primary site of \textbf{hematopoietic stem cells (HSCs)} (\textit{\autoref{core:hsc}}) and continuous blood cell generation in adult mammals, serving as the location for the complete maturation of all myeloid cells, NK cells, and B cells, as well as the site of the initial developmental stages of T cells.

%         The major anatomical sites containing active bone marrow include the femur, humerus, hip bones, and sternum. 
%       }{
%         \myfig{0.9}{immunology_figures/bone_marrow}{Bone marrow}{fig:bone_marrow}
%       }

%     \end{core}

%     \begin{core}{Thymus}

%       \sbs{0.55}{
%         The \textbf{thymus}, a small gland located above the heart, is the primary site of \textbf{T cell maturation} and development, involving a rigorous, stepwise selection process that tests for the correct functionality of T cells and their surface receptors. Cells that successfully survive this selection process enter systemic circulation, whereas the vast majority that fail undergo \textbf{apoptosis}.

%         Structurally, the thymus is organized into an \textbf{outer cortex} and an \textbf{inner medulla}, each supported by a distinct network of epithelial cells that guide developing \textbf{thymocytes} (\textit{immature, developing T cells}) through selection.
%       }{
%         \myfig{0.95}{immunology_figures/thymus}{Thymus}{fig:thymus}
%       }

%     \end{core}

%   \subsection{Secondary Immune System Organs}

%     \begin{core}{Secondary Immune System Organs}

%       The \textbf{secondary immune system organs} are the sites at which information regarding pathogens is coordinated, leading to the subsequent activation of immune cells.

%     \end{core}

%     \begin{core}{Lymphatic System (Fluids, Vessels, Nodes)}

%       \sbs{0.45}{
%         The \textbf{lymphatic system} functions as a \textbf{unidirectional} transport network, carrying fluid from peripheral tissues back toward the heart and ultimately the bloodstream.

%         \begin{itemize}

%           \item \textbf{Vessel \& Fluid}

%           \textbf{Lymphatic vessels} drain interstitial fluid, collecting antigens and leukocytes in the process; they facilitate the transport of immune cells but, unlike blood vessels, are devoid of \textbf{red blood cells (RBCs)}. Lymph fluid undergoes filtration as it passes through the lymph nodes.

%           \item \textbf{Lymph Nodes}

%             Lymph nodes serve primarily to [1] \textbf{trap incoming antigens} presented by APCs and [2] \textbf{provide a localized environment for lymphocytes}---T and B cells---to interact with these antigens, thereby generating a subsequent adaptive immune response.

%           The resulting activated lymphocytes exit the node through the efferent lymphatic vessels, ultimately entering the bloodstream via the thoracic duct and subclavian vein to traffic toward the site of infection.

%         \end{itemize}
%       }{
%         \myfig{1}{immunology_figures/lymph_system}{Lymphatic system}{fig:lymph_system}
%       }

%     \end{core}

%     \begin{core}{Spleen}

%       \sbs{0.45}{
%         The \textbf{spleen} functions strictly to \textbf{filter blood}, as opposed to lymph, and can be categorized into two tissue types.

%         \begin{itemize}

%           \item \textbf{Red Pulp} is responsible for recycling senescent RBCs and platelets from circulation via macrophage-mediated phagocytosis.

%           \item \textbf{White Pulp} is a crucial site for immune activation, organized into \textbf{periarteriolar lymphoid sheaths (PALS)} primarily containing T cells, and marginal \textbf{lymphoid follicles} primarily containing B cells.

%         \end{itemize}

%         \textit{Surgical removal of the spleen can significantly elevate an individual's susceptibility to bacterial infections.}
%       }{
%         \myfig{0.95}{immunology_figures/spleen}{Spleen}{fig:spleen}
%       }

%     \end{core}

%     \begin{core}{Mucosa Associated Lymphoid Tissues (MALT)}

%       \sbs{0.4}{
%         \textbf{Mucosa associated lymphoid tissues (MALT)} refer to the collection of lymphoid tissues situated across various mucosal surfaces, including the gastrointestinal (GALT), bronchial (BALT), and nasal (NALT) tracts; these mucosal systems represent the primary sites of pathogen entry into the body.

%         \textit{Highly organized MALT structures include the tonsils, the appendix, and \textbf{Peyer's patches} located in the intestinal tract.}

%         \textbf{Microfold (M) cells} (\textit{\autoref{fig:malt}}) are specialized, large epithelial cells located in areas of localized lymphoid concentration throughout the MALT, each possessing a basolateral pocket that houses several smaller immune cells.

%         The epithelial cells of the mucosa serve the following functions.
%       }{
%         \myfig{0.9}{immunology_figures/malt}{Mucosa associated lymphoid tissues; \textit{microfold cell present towards the left end of epithelial cell line}}{fig:malt}
%       }

%       \begin{itemize}

%         \item They facilitate the \textbf{transport of antigen samples} from the \textbf{lumen} (\textit{the space within vessels}) across the epithelial barrier via the specialized M cells.

%         \item They \textbf{secrete fluids} such as tears and mucus, which contain antimicrobial enzymes like lysozyme that degrade bacterial membranes.

%         \item They utilize ciliary action to \textbf{physically sweep away pathogens and debris}, a mechanism prominent in the respiratory and reproductive tracts.

%       \end{itemize}

%     \end{core}

%   \subsection{Tertiary Immune System Organs}

%   \begin{core}{Tertiary Immune System Organs}

%     A \textbf{tertiary immune organ} is an organized, non-encapsulated aggregate of immune and stromal cells that develops postnatally within non-lymphoid tissues at sites of persistent antigen exposure.

%   \end{core}

%   \begin{core}{Skin}

%     The \textbf{skin} is the largest organ of the human body and provides both innate and adaptive immune defense mechanisms. Its innate defense characteristics are as follows.

%     \begin{itemize}

%       \item \textbf{Keratinocytes}, the epithelial cells forming the outer epidermal layer, secrete \textbf{cytokines} to activate localized innate immune cells upon detecting a threat.

%       \item Upon cell death, keratinocytes leave behind a protective \textbf{physical barrier} composed of keratin intermediate filaments.

%       \item In the \textbf{epidermis}, sebaceous glands and hair follicles secrete \textbf{sebum}, antibacterial peptides, and \textbf{psoriasin} (\textit{a specific antimicrobial peptide}) to eliminate pathogenic bacteria while maintaining normal flora.

%     \end{itemize}

%     \sbs{0.5}{
%       \begin{itemize}

%         \item The underlying \textbf{dermis} consists of a deeper, thicker layer of mesenchymal and connective tissue that provides essential \textbf{blood supply and structural support}.

%       \end{itemize}

%       Conversely, the adaptive defense mechanisms of the skin are as follows.

%       \begin{itemize}

%         \item Keratinocytes, which are capable of expressing MHC~II receptors (\textit{\autoref{core:mhc_characteristics}}), and \textbf{Langerhans cells} (\textit{specialized dendritic cells of the skin}) present antigens to activate T$_\text{H}$ cells.

%         \item Additionally, the skin harbors \textbf{intra-epidermal lymphocytes}, a specialized population of resident T cells that continuously patrol the tissue.

%       \end{itemize}
%     }{
%       \myfig{0.95}{immunology_figures/skin}{Skin}{fig:skin}
%     }

%   \end{core}

% \section{Innate Immunity Processes}

%   \begin{core}{Innate vs Adaptive Immune Response}

%     The \textbf{innate immune response} is the body's first line of defense against pathogens. It is a rapid, \textbf{non-specific}, and evolutionary ancient defense mechanism that acts within minutes to hours of exposure. It is one that you can produce prior to exposure by a specific pathogen and doesn't require change to DNA/genes. All organisms have a form of innate immunity.

%     By contrast, the \textbf{adaptive response} is is the antigen-specific branch of the immune system that mounts a \textbf{tailored} defense against specific pathogen and creates long-lasting immunological memory. This response type takes days to develop during a primary infection because it requres the activation, gene rearrangement, clonal selection, and proliferation of lymphocytes with uniquely formed antigen receptors.

%   \end{core}

%   \subsection{Inflammation}

%     \begin{core}{Inflammation}

%       \textbf{Inflammation} represents a primary, coordinated physiological response that mobilizes systemic immune components to address and neutralize emerging pathogenic threats.

%     \end{core}

%     \begin{core}{Fever \& The Local Response}

%       \textbf{Fever} is a key characteristic of the inflammatory response, induced by the \textbf{hypothalamus}, which regulates the body's baseline temperature. It serves as an overt systemic indicator of immune up-regulation, a response historically correlated with improved survival rates during infection.

%       Any form of tissue injury, including internal trauma, inherently triggers a \textbf{localized inflammatory response}, characterized by the classical signs of \textbf{swelling} (\textit{due to vascular leakage}), \textbf{redness} (\textit{due to increased local blood flow}), \textbf{localized heat}, and \textbf{pain}. Collectively, these physiological changes serve to discourage behaviors that might exacerbate the initial injury.

%       \myfig{0.45}{immunology_figures/inflammation}{Localized inflammatory response; \textit{invading pathogens and physical tissue damage from a skin laceration elicit the classical signs of inflammation and the recruitment of immune cells to defend against infection}}{fig:inflammation}

%     \end{core}

%     \begin{core}{Immune Cell Extravasation}

%       \textbf{Extravasation} is the process by which immune cells exit the circulatory system, migrating across the endothelial lining of blood vessels and into surrounding tissue toward a site of infection, injury, or inflammation. This process is particularly relevant to \textbf{neutrophils}, which are largely confined to the vasculature prior to infection and are selectively recruited to tissues via extravasation.

%       \myfig{0.8}{immunology_figures/extravasation}{Extravastion of a neutrophil}{fig:extravasation}

%       \begin{enumerate}[label=(\arabic*)]

%         \item \textbf{Rolling}: Neutrophils form transient, low-affinity interactions with the endothelial surface, causing them to slowly roll along the vascular wall.

%         \item \textbf{Activation}: Endothelial and local cytokine signals induce conformational changes in neutrophil \textbf{integrins}.

%         \item \textbf{Adhesion}: High-affinity binding halts the rolling process, firmly arresting the cell on the endothelium.

%         \item \textbf{Transendothelial Migration}: Neutrophils move between endothelial cells to enter the inflamed tissue, directed by local cytokine gradients, where they will subsequently phagocytize pathogens.

%       \end{enumerate}

%     \end{core}

%   \subsection{Innate Pathogen Pattern Recognition}

%     \begin{core}[label=core:pamp]{Pathogen Associated Molecular Patterns (PAMPs)}

%       \textbf{Pathogen associated molecular patterns (PAMPs)} are conserved molecular structures shared broadly across classes of microorganisms and are not found on host cells, allowing the \textbf{innate} immune system to recognize them as generic markers of infection via dedicated \textbf{pattern recognition receptors} (\textit{\autoref{core:prr}}).

%       \begin{enumerate}

%         \item \textbf{Proteins} (\textit{**very few recognized innately})

%           \begin{itemize}

%             \item \ul{Flagellin}: The primary structural component of bacterial flagella (\textit{tails}).

%             \item \ul{Profilin}: A surface protein found in protozoa that facilitates their motility and host cell targeting.

%           \end{itemize}

%         \item \textbf{Carbs \& Glycopeptides}

%           \begin{itemize}

%             \item \ul{Zymosan}: A carbohydrate complex present in fungal cell walls.

%             \myfig{0.75}{immunology_figures/peptidoglycan}{Gram-positive (\textit{left}) and Gram-negative (\textit{right}) bacteria; \textit{peptidoglycan (burgandy) is present in both cell wall types}}{fig:peptidoglycan}
%             \vspace{-0.5cm}

%             \item \ul{Peptidoglycan}: A major structural polymer of both Gram-positive and Gram-negative bacterial cell walls shown in \autoref{fig:peptidoglycan}.

%           \end{itemize}

%         \item \textbf{Lipids}

%           \begin{itemize}

%             \item \ul{Lipopolysaccharide (LPS)}: Innately recognized lipids that are typically conjugated to carbohydrates or peptides, exemplified by the outer membrane component of Gram-negative bacteria (\textit{\autoref{fig:peptidoglycan}}).

%           \end{itemize}

%         \item \textbf{Nucleic Acids}

%           \begin{itemize}

%             \item \ul{Atypical Nucleic Acids}: Nucleic acids displaying configurations atypical of the host, such as single-stranded DNA or double-stranded RNA.

%           \end{itemize}

%       \end{enumerate}

%     \end{core}

%     \begin{core}[label=core:prr]{Pattern Recognition Receptors (PRRs)}

%       \textbf{Pattern recognition receptors (PRRs)} are host receptors that recognize PAMPs (\textit{\autoref{core:pamp}}), enabling the \textbf{innate} immune system to detect and respond to a broad range of microbial threats independent of prior antigen exposure.

%       \begin{enumerate}

%         \item \textbf{Extracellular PRRs}

%           \begin{itemize}

%             \item \ul{Lysozyme}: An enzyme that specifically \textbf{degrades the peptidoglycan} layer of bacterial cell walls (\textit{\autoref{fig:peptidoglycan}}).

%             \item \ul{Psoriasin}: A skin-specific \textbf{antimicrobial peptide (AMP)} that \textbf{disrupts bacterial cell membranes}; \textit{whereas AMPs are generally distributed systemically and share this membrane-disruptive mechanism, psoriasin is notable for its restriction to cutaneous tissue}.

%             \item \ul{Mannose-binding lectin (MBL)}: A plasma protein that binds microbial carbohydrates to initiate the \textbf{lectin pathway of the complement system} (\textit{AUTOREF}).

%             \item \ul{C-reactive protein (CRP)}: An acute-phase plasma protein that binds to specific surface molecules on \textbf{microbes} and \textbf{damaged host cells}.

%             \item \ul{Lipopolysaccharide-binding protein (LBP)}: A protein that \textbf{recognizes and binds lipopolysaccharide (LPS)} (\textit{\autoref{fig:peptidoglycan}}), the principal constituent of the Gram-negative bacterial outer membrane.

%           \end{itemize}

%         \item \textbf{Cytoplasmic PRRs}

%           \begin{itemize}

%             \item \ul{NOD proteins}: Intracellular sensors that detect degradation products of \textbf{bacterial cell walls} within the cytoplasm.

%           \end{itemize}

%         \item \textbf{Toll-like receptors (TLRs)} function as monomers or dimers and are distributed either \textbf{internally}, on endosomal or phagolysosomal membranes, or \textbf{externally}, on the plasma membrane facing the ECM.

%           \myfig{0.6}{immunology_figures/tlr}{Toll-like receptors; \textit{external (top left) TLRs sit on the plasma membrane facing the ECM and internal (bottom right) TLRs sit on the endosomal membrane}}{fig:tlr}
%           \vspace{-0.5cm}

%           Upon ligand binding, TLRs initiate signaling cascades that activate the transcription factor NF-κB, upregulating the \textbf{production of pro-inflammatory cytokines}.

%           \begin{table}[H]
%             \centering
%             \small
%             \begin{tabular}{@{}l l p{9cm}@{}}
%               \hline
%               \textbf{Category} & \textbf{TLR} & \textbf{Ligand Recognition} \\
%               \hline
%               \multirow{4}{*}{External}
%                 & TLR1, TLR2, TLR6 & Mediate the recognition of various \textbf{bacterial lipopeptides} and cell wall \textbf{lipoproteins}. \\
%                 & TLR4 & Specifically detects \textbf{LPS} from Gram-negative bacteria. \\
%                 & TLR5 & Specifically recognizes \textbf{bacterial flagellin}. \\
%                 & TLR10 & \textit{Unknown.} \\
%               \hline
%               \multirow{3}{*}{Internal}
%                 & TLR3 & Detects \textbf{viral double-stranded RNA}. \\
%                 & TLR7, TLR8 & Detect \textbf{single-stranded viral RNA sequences}. \\
%                 & TLR9 & Recognizes \textbf{unmethylated CpG dinucleotide motifs} commonly found in bacterial and viral DNA. \\
%               \hline
%               Non-Human & TLR11, TLR12, TLR13 & These receptors are absent in humans but are expressed in mice, functioning primarily to recognize \textbf{profilin} (\textit{an actin-binding protein}) and bacterial 23S rRNA. \\
%               \hline
%             \end{tabular}
%             \caption{Ligand specificity of TLRs, categorized by cellular localization.}
%             \label{tab:tlr_ligand}
%           \end{table}

%       \end{enumerate}

%     \end{core}

%   \subsection{Phagocytosis \& NOX Complexes}

%     \begin{core}{Phagocytosis}

%       \sbs{0.4}{
%         \textbf{Phagocytosis} refers to the engulfment of pathogens via the extension of \textbf{pseudopodia} around the target, leading to the formation a a \textbf{phagosome}. This process is triggered upstream by \textbf{PAMPs} (\textit{\autoref{core:pamp}}) and their recognition by \textbf{PRRs} (\textit{\autoref{core:prr}}) and primarily carried out by macrophages and neutrophils, which degrades the trapped pathogen via \textbf{NOX complexes} (\textit{\autoref{core:nox_complex}}).
%       }{
%         \myfig{1}{immunology_figures/phagocytosis}{Phagocytosis}{fig:phagocytosis}
%       }

%     \end{core}

%     \begin{core}[label=core:nox_complex]{NADPH Oxidase (NOX) Complexes}

%       \textbf{NADPH oxidase (NOX) complexes} are multi-subunit, membrane-bound enzymes whose primary biological function is the targeted generation of \textbf{reactive oxygen species (ROS)} through the transfer of electrons across biological membranes. 

%       \pgfig{
%         [
%           >=Stealth,
%           every node/.style={font=\sffamily},
%           membranefill/.style={fill=boxblue},
%           boxfill/.style={fill=boxblue, draw=black},
%         ]

%         % ===== Labels =====
%         \node[font=\sffamily\bfseries, anchor=west] (phagosome) at (0,11.8) {\underline{PHAGOSOME}};
%         \node[font=\sffamily\bfseries, anchor=west] (cytoplasm) at (0,0.4) {\underline{CYTOPLASM}};

%         % ===== Bacteria cloud =====
%         \node[cloud, draw, boxfill, aspect=2.2, minimum width=3.6cm, minimum height=2.2cm, cloud puffs=13] (bacteria) at (13.2,11.0) {Bacteria};

%         % ===== Degradation arrow =====
%         \node[font=\sffamily] at (8.6,12.15) {Degradation};

%         % ===== Membrane =====
%         \coordinate (memL) at (0.3,4.0);
%         \coordinate (memR) at (15.6,4.0);
%         \draw[boxfill] (0.3,4.6) rectangle (15.6,5.4);
%         \node[anchor=west] at (0.5,5.0) {\normalsize Phagosome Membrane};
%         \draw[boxfill] (0.8,4.0) rectangle (15.1,4.6);
%         \node[anchor=west] at (1.0,4.3) {\normalsize Microfilaments};

%         % ===== NOX Complex =====
%         \node[boxfill, rounded corners=8pt, minimum width=2.3cm, minimum height=3.4cm, align=center] (nox) at (6.0,4.9) {\Large NOX\\Complex};


%         % ===== Defensins ellipses =====
%         \node[boxfill, ellipse, minimum width=2.2cm, minimum height=1.1cm] (def1) at (11.0,1.6) {Defensins};
%         \node[boxfill, ellipse, minimum width=2.2cm, minimum height=1.1cm] (def2) at (10.3,0.2) {Defensins};
%         \node[boxfill, ellipse, minimum width=2.2cm, minimum height=1.1cm] (def3) at (12.9,0.75) {Defensins};

%         % ===== Chemical species labels =====
%         \node[font=\large] (O2) at (3.9,7.4) {\ce{2O2}};
%         \node[font=\large] (O2m) at (7.6,7.6) {\ce{2O2^-}};
%         \node[font=\large] (H2O2) at (7.0,10.2) {\ce{H2O2}};
%         \node[font=\large] (HOCl) at (3.0,10.1) {\ce{HOCl}};
%         \node[font=\large] (NADPH) at (3.9,2.45) {\ce{NADPH}};
%         \node[font=\large] (NADPp) at (7.8,2.45) {\ce{NADP+ + H+}};
%         \node[font=\large] (Htop) at (10.15,7.4) {\ce{H+}};
%         \node[font=\large] (Hbot) at (10.15,2.85) {\ce{H+}};
%         \node[font=\large] (Ktop) at (12.15,7.4) {\ce{K+}};
%         \node[font=\large] (Kbot) at (12.15,2.6) {\ce{K+}};

%         % ===== Arrows =====
%         % NADPH into NOX complex
%         \draw[->,thick] (NADPH.north east) to[bend left=45] (NADPp.north west);
%         % 2O2 -> 2O2-
%         \draw[->,thick] (O2.south east) to[bend right=45] (O2m.south west);
%         % 2O2- -> H2O2
%         \draw[->,thick] (O2m.north) -- (H2O2.south);
%         % H2O2 -> HOCl (elbow / stepped line)
%         \draw[->,thick] (H2O2.west) -- (HOCl.east);
%         % HOCl -> curved arrow up to Degradation arrow start
%         \draw[->,thick] (HOCl.north east) to[bend left=15] (10.7,11.75);
%         % NADP+ + H+ -> H+ (bottom)
%         \draw[->,thick] (NADPp.east) -- (Hbot.west);
%         % Vertical arrow through Proton channel: Hbot -> Htop
%         \draw[->,thick] (Hbot.north) -- (Htop.south);
%         % Vertical arrow through Kv channel: Kbot -> Ktop
%         \draw[->,thick] (Kbot.north) -- (Ktop.south);
%         % Defensins -> Bacteria (curved, right side, swings wide of the Kv channel arrow)
%         \draw[->,thick] ([xshift=0.3cm,yshift=0.2cm]def1.east) to[bend right=60] ([xshift=0.2cm,yshift=-0.2cm]bacteria.south east);
%         \draw[->,thick] (Htop.north) -- (H2O2.south);

%         % ===== Proton Channels (cylinder) =====
%         \begin{scope}
%           \coordinate (pcB) at (10.15,3.55);
%           \coordinate (pcT) at (10.15,5.65);
%           \fill[boxfill] (9.5,3.55) rectangle (10.8,5.65);
%           \draw[boxfill] (9.5,3.55) -- (9.5,5.65) (10.8,3.55) -- (10.8,5.65);
%           \node[align=center] at (10.15,4.55) {\small Proton\\\small Channels};
%         \end{scope}

%         % ===== Kv Channels (cylinder) =====
%         \begin{scope}
%           \coordinate (kcB) at (12.15,3.55);
%           \coordinate (kcT) at (12.15,5.65);
%           \fill[boxfill] (11.5,3.55) rectangle (12.8,5.65);
%           \draw[boxfill] (11.5,3.55) -- (11.5,5.65) (12.8,3.55) -- (12.8,5.65);
%           \node[align=center] at (12.15,4.55) {\small Kv\\\small Channels};
%         \end{scope}

%         \node[font=\bfseries\large] at (2,3.5) {(A)};
%         \node[font=\bfseries\large] at (2,10) {(B)};
%         \node[font=\bfseries\large] at (14,2) {(C)};
%         \node[font=\bfseries\large] at (13,7) {(D)};
%       }{NADPH oxidase (NOX) complex mechanism; \textit{(A) \textbf{Microfilaments package the exterior} of the vacuole, preventing swelling and maintaining the concentration of toxic compounds inside}; \textit{(B) Both \textbf{ROS} and \textbf{RNS} are made to attack the bacterial cell, this figure only shows the formation of ROS however for simplicity}; \textit{(C) \textbf{Defensin peptides (a type of AMP)} attack the pathogen by poking holes in its membrane, working best in a certain \ce{K+} concentration and tonicity}; \textit{(D) In order to \textbf{regulate the electrochemical gradient} made by transferring ions across the membrane, \ce{K+} is transferred into the phagosomal lumen by membrane channels}}{fig:phk}{0.8}
       
%     \end{core}

\section{Antibodies}

  \subsection{Immunoglobulin Structure \& Isotypes}

    \begin{core}{Immunoglobulin Superfamily}

      \sbs{0.55}{
        The \textbf{immunoglobulin superfamily} refers to a group of molecules containing at least one \textbf{immunoglobulin (Ig) domain}. Regardless of the specific protein it belongs to, almost all domains share a highly conserved 3D structure consisting of 70 to 110 amino acids folded into a compact lump of \textbf{β pleated sheets} that are pined together by a single \textbf{intradomain disulfide bond} between \textbf{two cysteine residues}

        It's natural confirmation results in the inner-facing side chains of the β sheets being primarily \textbf{hydrophobic}, driving them to pack itghly against each other, further stabilizing what is often referred to as the "\textit{sandwich}" structure.
      }{
        \myfig{0.65}{immunology_figures/ig_domain}{Three-dimensional structure of an immunoglobulin domain}{fig:ig_domain}
      }

    \end{core}

    \begin{core}[label=core:ig_receptor]{Immunoglobulin Receptors \& Antibodies}

      An \textbf{antibody}---also known simply as an \textbf{immunoglobulin} or \textbf{Ig}---is a highly specialized, Y-shaped protein composed of Ig domains that is produced by the immune system to recognize and neutralize foreign threats. An \textbf{Ig receptor} (\textit{or B cell receptor}) is an antibody bearing a membrane-spanning and cytosolic domain at its \textit{C-terminal} end. The general antibody structure can be divided into two \textbf{light (L)} and two \textbf{heavy (H)} chains.

      \begin{enumerate}

        \item \textbf{Light Chains}: The two light chains---$~25,000$ MW; $50,000$ MW total for both---are identical to one another and each consist of 2 Ig domains: one \textbf{variable} (\textit{\ce{V_L} in \autoref{fig:antibody}}) and one \textbf{constant} (\textit{\ce{C_L} in \autoref{fig:antibody}}).

          \myfig{0.55}{immunology_figures/light_chain}{Three-dimensional structure of an antibody light chain; \textit{variable Ig domain (left), joining segment (green), constant domain (blue)}}{fig:light_chain}

          \begin{itemize}

            \item \textbf{Constant Region (\ce{C_L})}: The carboxy-terminal (\ce{COOH}) portion of this region \textbf{links with the heavy chain}, contributing to the hinge structure that forms the base of the characteristic Y-shape. Two distinct isotypic variants of the light chain constant region exist: the \textbf{kappa (κ)} and \textbf{lambda (λ)} chains, the latter of which can be further divided into five structural subclasses.

            \item \textbf{Variable Region (\ce{V_L})}: This region contains three distinct loops extending from the terminus that together form the \textbf{hypervariable region}, constructed from three non-contiguous amino acid sequences and comprising approximately 15--20\% of the domain.

          \end{itemize}

        \pgfig{
          \tikzset{
            redbox/.style={rectangle, draw, fill={rgb,255:red,224;green,102;blue,102}, minimum width=3cm, minimum height=1cm, align=center, font=\bfseries},
            bluebox/.style={rectangle, draw, fill={rgb,255:red,61;green,133;blue,198}, text=white, minimum width=3cm, minimum height=1cm, align=center, font=\bfseries},
            lightbluebox/.style={rectangle, draw, fill={rgb,255:red,207;green,226;blue,243}, minimum width=3cm, minimum height=1cm, align=center, font=\bfseries},
          }

          % CENTER X
          \node[bluebox, rotate=45] at (0,0) {};
          \node[bluebox, rotate=-45] at (0,0) {\small Hinge-Bend Region};

          % RIGHT side
          \node[bluebox, rotate=45] at (2.2,2.2) {Constant Region};
          \node[lightbluebox, rotate=45] at (3.0,1.4) {\Large \ce{C_L}}; % 0.8 difference

          \node[redbox, rotate=45] at (4.4,4.4) {Variable Region};
          \node[redbox, rotate=45] at (5.2,3.6) {\Large \ce{V_L}}; % 0.8 difference

          \node[bluebox, rotate=-70] at (1.8,-2.8) {};
          \node[bluebox, rotate=-120] at (1.5,-6.0) {};

          % LEFT side
          \node[bluebox, rotate=-45] at (-2.2,2.2) {};
          \node[lightbluebox, rotate=-45] (lightbox) at (-3.0,1.4) {}; % 0.8 difference

          \node[redbox, rotate=-45] at (-4.4,4.4) {};
          \node[redbox, rotate=-45] at (-5.2,3.6) {}; % 0.8 difference

          \node[rectangle, draw, fill={rgb,255:red,182;green,215;blue,168}, minimum width=0.8cm, minimum height=0.2cm, align=center, font=\bfseries, rotate=-45] (dGroup) at (-5.45,4.9) {};

          \node[rectangle, draw, fill={rgb,255:red,254;green,228;blue,152}, minimum width=2.2cm, minimum height=0.2cm, align=center, font=\bfseries, rotate=-45] (jGroup) at (-4.39,3.84) {};

          \node[rectangle, draw, fill={rgb,255:red,254;green,228;blue,152}, minimum width=3cm, minimum height=0.2cm, align=center, font=\bfseries, rotate=-45] at (-4.93,3.87) {};

          \node[bluebox, rotate=70] at (-1.8,-2.8) {Constant Region};
          \node[bluebox, rotate=-60] at (-1.5,-6.0) {\small C-Term Domain};

          \node[cloud, draw, fill=boxblue, draw=black, aspect=2.2, minimum width=2.7cm, minimum height=2.2cm, cloud puffs=13] at (0,-3) {\small Oligosaccharide};

          \filldraw[fill=Red, draw=black] (0.4,0.4) circle (0.2);
          \filldraw[fill=Red, draw=black] (1.8,1) circle (0.2) node (linkageGroup) {};
          \filldraw[fill=Red, draw=black] (-1.8,1) circle (0.2);

          \filldraw[fill=Green, draw=black] (-4.4,4.6) circle (0.2);
          \filldraw[fill=Green, draw=black] (-4.6,4.4) circle (0.2);

          \filldraw[fill=Green, draw=black] (-5.2,3.8) circle (0.2);
          \filldraw[fill=Green, draw=black] (-5.4,3.6) circle (0.2) node (rGroup) {};

          \node[regular polygon, regular polygon sides=6, minimum size=2cm, draw=black, fill={rgb,255:red,142;green,124;blue,195}, thick, font=\bfseries, text=white] at (7,6.5) {\large Antigen};

          \draw[|-|, thick] (2.5,-0.2) -- (6.8,4.1) node[midway, below, sloped] {\bfseries \Large Light Chain};

          \draw[|-|, thick] plot[smooth] coordinates {(5,6) (-2,-1) (-3,-5) (-1.5,-7.6)};

          \node[rotate=45] at (2.3,4.0) {\bfseries \Large Heavy Chain};

          \node[rotate=45] at (-6.1,5.5) {\bfseries \Large CDR};

          \draw[|-|, thick] (-4,-7.6) -- (-4,-1) node[midway, above, xshift=-1.5cm] {\bfseries \Large Fc Stem};

          \draw[{latex}-, thick] (rGroup) -- ++(-1,-2) node[below, align=center] {Stabilizing\\Cysteine R\\Groups};

          \draw[{latex}-, thick] (linkageGroup) -- ++(3,-2) node[below right, align=center] {\bfseries \large (C)};

          \draw[{latex}-, thick] (dGroup) -- ++(3,2) node[right, align=center] {Diversity Group};

          \draw[{latex}-, thick] (jGroup) -- ++(3,2) node[right, align=center] {Joining Group};

          \draw[{latex}-, thick] (lightbox) -- ++(-3,-2) node[left, align=center] {Kappa (κ) OR\\Lambda (λ; \textit{5 subversions})};

          \node[] at (-7.5,5) {\bfseries \large (A)};
          \node[] at (0,-4.5) {\bfseries \large (B)};
        }{Antibody structure diagram; \textit{(A) \textbf{three loops} protrude from the ends of the variable domains and serve as the site of connection for a particular \textbf{antigen epitope}---each distinct CDR site is referred to as a distinct \textbf{idiotype}}; \textit{(B) an \textbf{oligosaccharide attaches to the second Ig domain}, determining how the antibody interacts with the rest of the immune system}; \textit{(C) both light--heavy and heavy--heavy linkages involve \textbf{covalent disulfide bonds} and strong \textbf{non-covalent forces}}}{fig:antibody}{0.8}

        \item \textbf{Heavy Domain}: The two heavy chains---$~50,000$ MW; $100,000$ MW for both---are also identical to one another and each consist of 4 or 5 Ig domains: one variable domain and the remaining constant domains, which serve the same general function across distinct antibodies.

          \begin{itemize}

            \item \textbf{Variable Region ($V_H$)}: Like the \ce{V_L} domain, three terminal loops form the \textbf{hypervariable region} of the heavy chain. When spatially paired with the three corresponding loops of the light chain \ce{V_L}, these loops collectively form the \textbf{complementarity-determining regions (CDRs)}, which are responsible for unique, hyperspecific antigen binding.

            \item \textbf{First Constant Domain}

            \item \textbf{Hinge-Bend Region}: This region is either a flexible \textbf{hinge} or a rigid \textbf{bend}, depending on the specific \textbf{Ig isotype} (\textit{\autoref{core:isotype}}). In IgM and IgE (\textit{possessing μ and ε heavy chains, respectively}), this region consists of additional, highly rigid immunoglobulin domains. In IgG, IgA, and IgD (\textit{possessing γ, α, and δ heavy chains}), it instead consists of a proline-rich amino acid sequence stabilized by disulfide bonds, which provides substantial structural flexibility.

            \item \textbf{Next Constant Domain}: This domain serves as the primary site for post-translational \textbf{oligosaccharide attachment}, which spatially opens the domains to the aqueous environment, which is essential for facilitating subsequent interactions.

            \item \textbf{Carboxy-Terminal (C-term) Domain}: This terminal domain dictates the antibody's final localization: whether it is retained as a \textbf{membrane-bound receptor} (\textit{on naive or memory B cells}) or \textbf{secreted extracellularly} (\textit{by \textbf{plasma cells}}). The secreted variant terminates in a short, exclusively hydrophilic amino acid sequence, whereas the membrane-bound variant terminates in a specialized sequence containing a hydrophobic transmembrane span followed by a short, cytoplasmic hydrophilic tail.

          \end{itemize}

      \end{enumerate}

    \end{core}

    \begin{core}[label=core:isotype]{Immunoglobulin Isotypes (\textit{H Chain Differences})}

      There are five classes of antibodies---termed IgM, IgD, IgG, IgA, and IgE---which differ in function as well as in the exact amino acid sequence and conformation of the Fc stem. In all cases, the basic unit consists of two heavy chains and two light chains, and the light chain may be either κ or λ, as discussed in \autoref{core:ig_receptor}. The classification of antibodies into these five classes depends on the structural elements of and functions dependent on the heavy chain as outlined below.

      \begin{itemize}

        \item \textbf{Flexible hinge/Rigid bend}: The number of constant (C) Ig domains present, which determines whether the heavy chain possesses a flexible hinge region or a more rigid bend.

        \item \textbf{Oligocsaccharide}: A small carbohydrate group is added to the second domain from the C-terminus, with the specific carbohydrate varying according to antibody type.

        \item \textbf{J Chain}: The presence of a J chain allows certain antibodies to form compound structures capable of crossing epithelial layers.

        \item \textbf{Subclasses}: Distinct subcategories exist within a given antibody class (\textit{e.g., IgA1, IgA2, IgA3, IgA4}).

      \end{itemize}

      \begin{table}[H]
        \centering
        \begin{tabular}{
          >{\itshape}c
          >{\centering\arraybackslash}p{0.5cm}
          >{\centering\arraybackslash}p{1.2cm}
          >{\arraybackslash}p{0.5cm}
          >{\raggedright\arraybackslash}p{3cm}
          >{\raggedright\arraybackslash}p{2.5cm}
        }
          \hline
          \textbf{IgX} & \textbf{Hinge/ Bend} & \textbf{J chain} & \textbf{Sub- classes} & \textbf{Timing} & \textbf{Membrane Spanning} \\
          \hline
          M & Rigid & \checkmark & --- & First class produced in maturing B cells & Na\"ive \& memory \\
          D & Hinge & $\times$ & --- & Produced as Ig receptor on mature, na\"ive B cells & Na\"ive cells (\textit{rarely soluble/memory}) \\
          G & Hinge & $\times$ & 4 & After class switching in activated B cells & Memory cells \\
          A & Hinge & \checkmark & 2 & After class switching in activated B cells & Memory cells \\
          E & Bend & $\times$ & 1 & After class switching in activated B cells & Memory cells \\
          \hline
        \end{tabular}
        \caption{Structural and functional properties distinguishing the five immunoglobulin classes.}
        \label{tab:ig_classes}
      \end{table}

      The following is a more in depth look into the functional roles and characteristics of each Ig subclass.

      \begin{enumerate}

        \item \textbf{IgM} (\textit{Memory \& Naive B Cells})

          IgM is structurally defined by a \textbf{rigid bend} region comprising four constant domains and serves as a critical first-line defense, representing the first antibody isotype secreted into plasma during a \textbf{primary immune response}. 

          \myfig{0.45}{immunology_figures/igm_pentamer}{IgM Pentamer}{fig:igm_pentamer}

          When expressed as a membrane-bound B cell receptor, it exists strictly as a monomer; however, its secreted soluble form assembles into a large \textbf{pentameric structure} possessing ten distinct antigen-binding sites. This pentamer is stabilized by a central \textbf{J chain}, which both aids in \textbf{mucosal secretion} and, through its multimeric geometry, renders IgM highly efficient at \textbf{cross-linking antigens} and initiating the \textbf{complement cascade}. 

          As with IgG, this potent complement activation drives rapid \textbf{opsonization} of targets and promotes their subsequent phagocytic destruction. Despite its larger size, IgM can also undergo \textbf{transcytosis} across mucosal barriers.

        \item \textbf{IgD} (\textit{Naive B Cells})

          IgD possesses a \textbf{flexible hinge} structure and assists in the \textbf{initial recognition of antigen by naive B cells}, predominantly operating as a \textbf{membrane-bound B cell receptor} rather than as a secreted antibody; accordingly, it is detected in plasma only at trace levels, accounting for roughly 0.2\% of total serum immunoglobulins.

        \item \textbf{IgG} (\textit{Memory B Cells})

          \myfig{0.7}{immunology_figures/igg_subclasses}{IgG subclasses (\textit{IgG1, IgG2, IgG3, and IgG4})}{fig:igg_subclasses}

          IgG is characterized by a \textbf{flexible hinge} structure comprising three constant domains and represents the \textbf{predominant secreted antibody class}, responsible for widespread systemic defense against both bacterial and viral pathogens. It is highly effective at activating the \textbf{complement system}, a process that leads to pathogen \textbf{opsonization} and consequently enhances phagocytic clearance. IgG further diversifies into four subclasses, numbered 1 through 4, which vary in their biological specificity and effector function, as summarized below.

          \begin{table}[H]
            \centering
            \begin{tabular}{
              >{\itshape}c
              >{\raggedright\arraybackslash}p{3.5cm}
              >{\raggedright\arraybackslash}p{3.5cm}
            }
              \hline
              \textbf{Subclass} & \textbf{Complement Activation} & \textbf{Fc Receptor Binding} \\
              \hline
              IgG1 & Moderately activates & Binds tightly \\
              IgG2 & Weakly activates & Binds weakly \\
              IgG3 & Strongly activates & Binds tightly (high disulfide content) \\
              IgG4 & Does not activate & Binds weakly \\
              \hline
            \end{tabular}
            \caption{Comparative complement activation and Fc receptor binding across IgG subclasses.}
            \label{tab:igg_subclasses}
          \end{table}

        \item \textbf{IgA} (\textit{Memory B Cells})

          \myfig{0.6}{immunology_figures/iga_dimer}{IgA dimer}{fig:iga_dimer}

          IgA possesses a \textbf{flexible hinge} structure and exists in two subclasses, IgA1 and IgA2. Its primary role involves \textbf{immune surveillance} and \textbf{defense at mucosal and epithelial barriers}, and it is secreted in substantial quantities---up to 15 grams per day---into mucus, tears, saliva, and breast milk. While IgA circulates in the serum \textbf{largely as a monomer}, it is \textbf{assembled and secreted across epithelia primarily as a dimer}, covalently linked by a \textbf{J chain} that facilitates its secretion into \textbf{mucus layers} and localizes the antibody to various \textbf{epithelial linings}; the active transport mechanism responsible for moving these antibodies across epithelial cell layers is termed \textbf{transcytosis}.

        \item \textbf{IgE} (\textit{Memory B Cells})

          IgE possesses a \textbf{rigid bend} region owing to an additional constant domain, and as a monomer it superficially resembles IgM, though it differs somewhat in amino acid sequence. Physiologically, it is adapted for immunological defense against \textbf{large parasitic helminths, or worms}. IgE is also deeply implicated in allergic pathogenesis, as its constant regions bind with exceptionally \textbf{high affinity to specialized Fc receptors} on the surface of mast cells and basophils, triggering the \textbf{characteristic allergic response}.

      \end{enumerate}

    \end{core}

  \subsection{Immunoglobulin Expression}

    \begin{core}{Immunoglobulin DNA to Expression}

      <FIGURE>

    \end{core}

    \begin{core}{Immunoglobulins in Action}

      Antigen recognition occurs at the six loops located at the tips of the antibody's Y-shaped arms, mediated by the same weak, non-covalent interactions responsible for enzyme-substrate binding. As in enzyme-substrate interactions, this binding process can involve induced fit, in which both the antibody and the antigen undergo conformational distortion to achieve optimal complementarity.

      Antigens are usually proteins, and several structural features govern how B-cell epitopes are recognized. Epitopes are typically located at the surface of a protein, often corresponding to regions that protrude outward and are thus most accessible to antibody binding. The protein's tertiary structure, or three-dimensional shape, is critical to this recognition; denatured proteins consequently fail to bind, since the epitope's native conformation has been lost. Quaternary structure can also play a role, as an antibody may recognize an epitope formed at the junction between two distinct associated peptides. Because of this dependence on three-dimensional conformation, B-cell epitopes may be formed either by contiguous (sequential) stretches of amino acids or by residues that are non-sequential in the primary structure but are brought into spatial proximity through protein folding.

    \end{core}

  \subsection{Immunoglobulin Genetics}

    % Since we are diploid, there is a pair of each chromosome and therefore each of the following arrangements may occur twice in the genome
    Because somatic cells are diploid, they possess two alleles for every relevant chromosome. Consequently, each genetic rearrangement event has the potential to occur on either allele within the genome.
    % Many of the elements of these genes are repeated variations of interchangeable parts
    The genetic loci encoding immunoglobulins consist largely of tandemly repeated, homologous segments that function as interchangeable coding blocks during somatic recombination.

    \begin{core}{L Chain Gene Sequences}

      \begin{enumerate}

        \item \textbf{λ L Chain}

          \pgfig{
            \hbandreset{0.2} % starting x-position
            \draw[dashed] (0,0) node[left] {5'} -- (16,0) node[right] {3'};

            \draw[hband black={0.2}{L}];
            \hbandshift{0.1}
            \draw[hband pink={0.8}{V}];
            \hbandshift{0.4}

            \draw[hband black={0.2}{L}];
            \hbandshift{0.1}
            \draw[hband pink={0.8}{V}];
            \hbandshift{1.0}

            \draw[hband black={0.2}{L}];
            \hbandshift{0.1}
            \draw[hband pink={0.8}{V}];
            \hbandshift{0.4}

            \draw[hband yellow={0.3}{J}];
            \hbandshift{0.2}
            \draw[hband blue={1.2}{C}];
            \hbandshift{0.4}

            \draw[hband yellow={0.3}{J}];
            \hbandshift{0.2}
            \draw[hband blue={1.2}{C}];
            \hbandshift{0.4}

            \draw[hband yellow={0.3}{J}];
            \hbandshift{0.2}
            \draw[hband blue={1.2}{C}];
            \hbandshift{0.4}

            \draw[hband yellow={0.3}{J}];
            \hbandshift{0.2}
            \draw[hband blue={1.2}{C}];
            \hbandshift{0.4}

            \draw[hband yellow={0.3}{J}];
            \hbandshift{0.2}
            \draw[hband blue={1.2}{C}];
            \hbandshift{0.4}
          }{caption}{fig:ph1}{0.9}

          % Chromosome 22
          Localized to Chromosome 22.
          % V Regions: Approximately 30-40
          V Segments: Approximately 30 to 40 distinct coding blocks.
          % J Regions: 4
          J Segments: 4 distinct coding blocks.
          % C Regions: 5 (2 functional [Cλ1, Cλ2], 3 nonfunctional [Cλ3, Cλ4, Cλ5])
          C Segments: 5 blocks total, containing 2 functional exons (Cλ1, Cλ2) and 3 pseudogenes (Cλ3, Cλ4, Cλ5).

        \item \textbf{κ L Chain}

          \pgfig{
            \hbandreset{0.2}
            \draw[dashed] (0,0) node[left] {5'} -- (13,0) node[right] {3'};

            \draw[hband black={0.2}{L}];
            \hbandshift{0.1}
            \draw[hband pink={0.8}{V}];
            \hbandshift{0.4}

            \draw[hband black={0.2}{L}];
            \hbandshift{0.1}
            \draw[hband pink={0.8}{V}];
            \hbandshift{1.0}

            \draw[hband black={0.2}{L}];
            \hbandshift{0.1}
            \draw[hband pink={0.8}{V}];
            \hbandshift{0.6}

            \draw[hband yellow={0.3}{J}];
            \hbandshift{0.8}

            \draw[hband yellow={0.3}{J}];
            \hbandshift{0.8}

            \draw[hband yellow={0.3}{J}];
            \hbandshift{0.8}

            \draw[hband yellow={0.3}{J}];
            \hbandshift{0.8}

            \draw[hband yellow={0.3}{J}];
            \hbandshift{0.8}

            \draw[hband blue={1.2}{C}];
          }{caption}{fig:ph2}{0.9}

          % Chromosome 2
          Localized to Chromosome 2.
          % V Regions: Approximately 35-40
          V Segments: Approximately 35 to 40 distinct coding blocks.
          % J Regions: 5
          J Segments: 5 distinct coding blocks.
          % C Regions: 1 (Cκ)
          C Segments: 1 functional coding block (Cκ).

      \end{enumerate}

      <FIGURE details>

      % content about the details
      Further molecular details concerning these sequence loci will be elaborated upon below.

    \end{core} 

    \begin{core}{H Chain Gene Sequences}

      \pgfig{
        \hbandreset{0.2}
        \draw[dashed] (0,0) node[left] {5'} -- (16,0);

        \draw[hband black={0.2}{L}];
        \hbandshift{0.1}
        \draw[hband pink={0.8}{V}];
        \hbandshift{1.0}

        \draw[hband black={0.2}{L}];
        \hbandshift{0.1}
        \draw[hband pink={0.8}{V}];
        \hbandshift{0.6}

        \draw[hband green={0.2}{D}];
        \hbandshift{0.8}

        \draw[hband green={0.2}{D}];
        \hbandshift{0.6}

        \draw[hband yellow={0.3}{J}];
        \hbandshift{0.8}

        \draw[hband yellow={0.3}{J}];
        \hbandshift{0.8}

        \draw[hband yellow={0.3}{J}];
        \hbandshift{0.8}

        \draw[hband yellow={0.3}{J}];
        \hbandshift{0.8}

        \draw[hband yellow={0.3}{J}];
        \hbandshift{0.8}

        \hbandreset{0.2}

        \draw[dashed] (0,-2) -- (16,-2) node[right] {3'};

        \begin{scope}[shift={(0,-2)}]
          \draw[hband blue={1.2}{C}];
          \hbandshift{0.5}

          \draw[hband blue={1.2}{C}];
          \hbandshift{0.5}

          \draw[hband blue={1.2}{C}];
          \hbandshift{0.5}

          \draw[hband blue={1.2}{C}];
          \hbandshift{0.5}

          \draw[hband blue={1.2}{C}];
          \hbandshift{0.5}

          \draw[hband blue={1.2}{C}];
          \hbandshift{0.5}

          \draw[hband blue={1.2}{C}];
          \hbandshift{0.5}

          \draw[hband blue={1.2}{C}];
          \hbandshift{0.5}

          \draw[hband blue={1.2}{C}];
          \hbandshift{0.5}
        \end{scope}

      }{caption}{fig:ph3}{0.85}

      % Chromosome 14
      Localized to Chromosome 14.

      % V Regions: Approximately 40
      V Segments: Approximately 40 functional segments.

      % D Regions: About 23
      D Segments: Approximately 23 functional segments.

      % J Regions: Approximately 6
      J Segments: Approximately 6 functional segments.

      % C Regions: 8 (corresponding to the different antibody isotypes; the Cμ and Cδ count as one C region)
      C Segments: 8 distinct domains that correspond directly to the various antibody isotypes (noting that the closely linked Cμ and Cδ segments are functionally transcribed together prior to alternative splicing).

      % [Cμ/Cδ] [Cγ3] [Cγ1] [Cα1] [Cγ2] [Cγ4] [Cε] [Cα2] in order**
      The specific downstream genomic arrangement of these constant regions follows the order: Cμ/Cδ, Cγ3, Cγ1, Cα1, Cγ2, Cγ4, Cε, and Cα2.

      <FIGURE details>

      % If we look deeper into the regions that code for the heavy chain, we can see more specific variety of DNA sequences
      Examining the heavy chain constant regions at a higher resolution reveals an organized substructure of coding sequences:

        % 3-4 distinguished DNA sequences that represent the different Ig domains among Ig isotypes
        Each region contains 3 to 4 independent exons that correspond to the individual globular domains comprising the structural heavy chain of that isotype.

        % Two downstream regions after the Ig domain sequences represent the membrane spanning extensions
        Situated downstream of the primary structural exons are two highly conserved sequences coding for the hydrophobic membrane-spanning extensions.

      % ***Both these regions are followed by a place where you can place a poly-A-tail to decide if the antibody will be soluble or membrane bound
      Importantly, separate polyadenylation (poly-A) signal sites exist after both the secreted-sequence exons and the membrane-spanning exons, dictating through alternative RNA processing whether the final synthesized antibody will be secreted or membrane-anchored.

    \end{core}

    \subsubsection{L/H Chain Gene Expression}

      \begin{core}{λ L Chain Gene Expression}

        \begin{enumerate}[label=(\arabic*)]

          \item \textbf{REARRANGEMENT} [Rearranged Gene] --- \textit{Nucleus}

            \pgfig{
              \hbandreset{0.2} % starting x-position
              \draw[dashed] (0,0) node[left] {5'} -- (7,0) node[right] {3'};

              \draw[hband black={0.2}{L}];
              \hbandshift{0.1}
              \draw[hband pink={0.8}{V}];
              \hbandshift{0.4}

              \draw[hband black={0.2}{L}];
              \hbandshift{0.1}
              \draw[hband pink={0.8}{V}];
              \draw[hband yellow={0.3}{J}];
              \hbandshift{0.2}
              \draw[hband blue={1.2}{C}];
              \hbandshift{0.4}

              \draw[hband yellow={0.3}{J}];
              \hbandshift{0.2}
              \draw[hband blue={1.2}{C}];
              \hbandshift{0.4}

              \hbandreset{0.2} % starting x-position
              \begin{scope}[shift={(0,-2)}]
                \hbandreset{0.2} % starting x-position
                \draw[dashed] (0,0) node[left] {5'} -- (16,0) node[right] {3'};

                \draw[hband black={0.2}{L}];
                \hbandshift{0.1}
                \draw[hband pink={0.8}{V}];
                \hbandshift{0.4}

                \draw[hband black={0.2}{L}];
                \hbandshift{0.1}
                \draw[hband pink={0.8}{V}];
                \hbandshift{1.0}

                \draw[hband black={0.2}{L}];
                \hbandshift{0.1}
                \draw[hband pink={0.8}{V}];
                \draw[hband yellow={0.3}{J}];
                \hbandshift{0.2}
                \draw[hband blue={1.2}{C}];
                \hbandshift{0.4}

                \draw[hband yellow={0.3}{J}];
                \hbandshift{0.2}
                \draw[hband blue={1.2}{C}];
                \hbandshift{0.4}

                \draw[hband yellow={0.3}{J}];
                \hbandshift{0.2}
                \draw[hband blue={1.2}{C}];
                \hbandshift{0.4}

                \draw[hband yellow={0.3}{J}];
                \hbandshift{0.2}
                \draw[hband blue={1.2}{C}];
                \hbandshift{0.4}

                \draw[hband yellow={0.3}{J}];
                \hbandshift{0.2}
                \draw[hband blue={1.2}{C}];
                \hbandshift{0.4}
              \end{scope}

              \hbandreset{0.2} % starting x-position
              \begin{scope}[shift={(0,-4)}]
                \draw[dashed] (0,0) node[left] {5'} -- (4,0) node[right] {3'};

                \draw[hband black={0.2}{L}];
                \hbandshift{0.1}
                \draw[hband pink={0.8}{V}];
                \draw[hband yellow={0.3}{J}];
                \hbandshift{0.2}
                \draw[hband blue={1.2}{C}];
                \hbandshift{0.4}

              \end{scope}

            }{caption}{fig:ph4}{0.9}

            % 1 LV and 1 JC pair come together at random, clipping out any JC pairs between the selected VJC combination and activating the promoter of the selected VL only
            A single variable (V) segment undergoes somatic recombination with a single joining (J) segment. This DNA rearrangement physically excises any intervening genetic sequences and functionally activates only the promoter immediately upstream of the successfully joined V segment.

            \pgfig{
              \hbandreset{0.2} % starting x-position
              \draw[dashed] (0,0) node[left] {5'} -- (9,0) node[right] {3'};

              \draw[hband black={0.2}{L}];
              \hbandshift{0.1}
              \draw[hband pink={0.8}{V}];
              \hbandshift{5}
              \draw[hband yellow={0.3}{J}];
              \hbandshift{0.2}
              \draw[hband blue={1.2}{C}];
              \hbandshift{0.4}

              \draw[fill=imgray, draw=none] (1.3,-0.3) rectangle (2.0,0.3);
              \draw[fill=impurple, draw=none] (2.0,-0.3) rectangle (2.7,0.3);
              \draw[fill=imorange, draw=none] (3.3,-0.3) -- (2.7,-0.3) -- (2.7,0.3) -- cycle;
              \draw (2.7,0.3) -- (1.3,0.3) -- (1.3,-0.3) -- (3.3,-0.3) -- cycle; 

              \draw[fill=imlightpurple, draw=none] (4.9,-0.3) rectangle (5.6,0.3);
              \draw[fill=imgray, draw=none] (5.6,-0.3) rectangle (6.3,0.3);
              \draw[fill=imorange, draw=none] (4.3,0.3) -- (4.9,0.3) -- (4.9,-0.3) -- cycle;
              \draw (4.3,0.3) -- (6.3,0.3) -- (6.3,-0.3) -- (4.9,-0.3) -- cycle; 

            }{caption}{fig:rss1}{0.9}

            % reference to the mechanism core concept
            This process relies on the Recombination Signal Sequence (RSS) mechanism detailed in a subsequent section.

          \item \textbf{TRANSCRIPTION} [Primary Transcript] --- \textit{Nucleus}

            \pgfig{
              \hbandreset{0.2} % starting x-position
              \draw[dashed] (0,0) node[left] {5'} -- (4,0) node[right] {3'};

              \draw[hband black={0.2}{L}];
              \hbandshift{0.1}
              \draw[hband pink={0.8}{V}];
              \draw[hband yellow={0.3}{J}];
              \hbandshift{0.2}
              \draw[hband blue={1.2}{C}];
              \hbandshift{0.4}

            }{caption}{fig:ph5}{0.9}

            % RNA polymerase transcribes the VL and JC into a primary transcript, stopping at the stop signal present at the end of the C region sequence
            RNA polymerase generates a primary transcript covering the rearranged VJ sequence, continuing through the downstream constant (C) region until it terminates at the designated transcription stop signal.

          \item \textbf{PROCESSING} [Processed Message] --- \textit{Nucleus}

            \pgfig{
              \hbandreset{0.2} % starting x-position
              \draw[dashed] (0,0) node[left] {5'} -- (3,0) node[right] {AA(A\ce{_n})AA 3'};

              \draw[hband black={0.2}{L}];
              \draw[hband pink={0.8}{V}];
              \draw[hband yellow={0.3}{J}];
              \draw[hband blue={1.2}{C}];

            }{caption}{fig:ph6}{0.9}

            % Enzymes removes the introns from between the VL and the JC
            Spliceosomes process the primary RNA, removing non-coding intronic sequences located between the rearranged VJ unit and the C domain.

            % Adds a poly-A-tail (makes a RNA more stable)
            Post-transcriptional modification includes the addition of a poly-A tail at the 3' end to ensure messenger RNA stability.

          \item \textbf{TRANSLATION} [Nascent Peptide] --- \textit{Into the RER}

            \pgfig{
              \hbandreset{0.2} % starting x-position
              \draw[dashed] (0,0) node[left] {N} -- (3,0) node[right] {C};

              \draw[hband black={0.2}{L}];
              \draw[hband pink={0.8}{V}];
              \draw[hband yellow={0.3}{J}];
              \draw[hband blue={1.2}{C}];

            }{caption}{fig:ph7}{0.9}

            % The ribosome attaches to the message, begins translating, attaches to the Rough Endoplasmic Reticulum (RER), and pushes the nascent peptide into the ER lumen
            Ribosomal complexes bind the mature mRNA in the cytosol. Upon synthesizing the initial targeting sequence, the ribosome docks to the Rough Endoplasmic Reticulum (RER), allowing the nascent polypeptide chain to be translocated directly into the ER lumen.

            % The initial L sequence binds factors that cause the ribosome to attach to the RER
            This docking is mediated by a hydrophobic leader (L) sequence at the N-terminus of the developing protein, which interacts with signal recognition particles.

          \item \textbf{L REGION SNIPPING} [Mature Chain] --- \textit{ER Lumen}

            \pgfig{
              \hbandreset{0.2} % starting x-position
              \draw[dashed] (0,0) node[left] {N} -- (2.8,0) node[right] {C};

              \draw[hband pink={0.8}{V}];
              \draw[hband yellow={0.3}{J}];
              \draw[hband blue={1.2}{C}];

            }{caption}{fig:ph8}{0.9}

            % Enzymes clip off the leader, leaving a chain with a V(D)J domain and a constant domain
            Signal peptidases within the ER lumen cleave the N-terminal leader sequence, producing a mature light chain polypeptide consisting solely of the variable (VJ) domain and the constant (C) domain.

        \end{enumerate} 

      \end{core}

      \begin{core}{κ L Chain Gene Expression}

        \begin{enumerate}[label=(\arabic*)]

          \item \textbf{REARRANGEMENT} [Rearranged Gene] --- \textit{Nucleus}

            \pgfig{
              \hbandreset{0.2}
              \draw[dashed] (0,0) node[left] {5'} -- (9,0) node[right] {3'};

              \draw[hband black={0.2}{L}];
              \hbandshift{0.1}
              \draw[hband pink={0.8}{V}];
              \hbandshift{0.4}

              \draw[hband black={0.2}{L}];
              \hbandshift{0.1}
              \draw[hband pink={0.8}{V}];
              \draw[hband yellow={0.3}{J}];
              \hbandshift{0.8}

              \draw[hband yellow={0.3}{J}];
              \hbandshift{0.8}

              \draw[hband yellow={0.3}{J}];
              \hbandshift{0.8}

              \draw[hband yellow={0.3}{J}];
              \hbandshift{0.8}

              \draw[hband blue={1.2}{C}];
            }{caption}{fig:ph9}{0.9}

            % 1 VL and 1 J gene region come together, clipping out any J regions between the selected LVJ combination and again activating only the promoter of the selected VL only
            Similar to the lambda chain, a single V segment undergoes recombination with a single J segment. Any intervening J sequences are permanently excised from the DNA, and the upstream promoter corresponding exclusively to the chosen V segment becomes functionally active.

            \pgfig{
              \hbandreset{0.2} % starting x-position
              \draw[dashed] (0,0) node[left] {5'} -- (9,0) node[right] {3'};

              \draw[hband black={0.2}{L}];
              \hbandshift{0.1}
              \draw[hband pink={0.8}{V}];
              \hbandshift{5}
              \draw[hband yellow={0.3}{J}];
              \hbandshift{0.2}
              \draw[hband blue={1.2}{C}];
              \hbandshift{0.4}

              \draw[fill=imgray, draw=none] (1.3,-0.3) rectangle (2.0,0.3);
              \draw[fill=imlightpurple, draw=none] (2.0,-0.3) rectangle (2.7,0.3);
              \draw[fill=imorange, draw=none] (2.7,-0.3) -- (2.7,0.3) -- (3.3,0.3) -- cycle;
              \draw (3.3,0.3) -- (1.3,0.3) -- (1.3,-0.3) -- (2.7,-0.3) -- cycle; 

              \draw[fill=impurple, draw=none] (4.9,-0.3) rectangle (5.6,0.3);
              \draw[fill=imgray, draw=none] (5.6,-0.3) rectangle (6.3,0.3);
              \draw[fill=imorange, draw=none] (4.3,-0.3) -- (4.9,-0.3) -- (4.9,0.3) -- cycle;
              \draw (4.9,0.3) -- (6.3,0.3) -- (6.3,-0.3) -- (4.3,-0.3) -- cycle; 

            }{caption}{fig:rss2}{0.9}

            % reference to RSS mechanism core concept
            This DNA rearrangement is mediated by the RSS molecular mechanism.

          \item \textbf{TRANSCRIPTION} [Primary Transcript] --- \textit{Nucleus}

            \pgfig{
              \hbandreset{0.2}
              \draw[dashed] (0,0) node[left] {5'} -- (7.5,0) node[right] {3'};

              \draw[hband black={0.2}{L}];
              \hbandshift{0.1}
              \draw[hband pink={0.8}{V}];
              \draw[hband yellow={0.3}{J}];
              \hbandshift{0.8}

              \draw[hband yellow={0.3}{J}];
              \hbandshift{0.8}

              \draw[hband yellow={0.3}{J}];
              \hbandshift{0.8}

              \draw[hband yellow={0.3}{J}];
              \hbandshift{0.8}

              \draw[hband blue={1.2}{C}];
            }{caption}{fig:ph10}{0.9}

            % RNA polymerase transcribes the beginning Vl, the selected J (and all the downstream Js), and the constant region
            Transcription initiates at the active promoter, generating a primary RNA transcript that spans the chosen V segment, the selected J segment, any remaining downstream unselected J segments, and the single Cκ constant region.

          \item \textbf{PROCESSING} [Processed Message] --- \textit{Nucleus}

            \pgfig{
              \hbandreset{0.2} % starting x-position
              \draw[dashed] (0,0) node[left] {5'} -- (3,0) node[right] {AA(A\ce{_n})AA 3'};

              \draw[hband black={0.2}{L}];
              \draw[hband pink={0.8}{V}];
              \draw[hband yellow={0.3}{J}];
              \draw[hband blue={1.2}{C}];

            }{caption}{fig:kl1}{0.9}

            % Any introns and extra Js get clipped out
            RNA splicing enzymes excise all intronic sequences and the unselected downstream J segments, yielding a contiguous coding message.

            % Adds a poly-A-tail (makes a RNA more stable)
            A poly-A tail is appended to stabilize the transcript for cytosolic export.

          \item \textbf{TRANSLATION} [Nascent Peptide] --- \textit{Into the RER}

            \pgfig{
              \hbandreset{0.2} % starting x-position
              \draw[dashed] (0,0) node[left] {N} -- (3,0) node[right] {C};

              \draw[hband black={0.2}{L}];
              \draw[hband pink={0.8}{V}];
              \draw[hband yellow={0.3}{J}];
              \draw[hband blue={1.2}{C}];

            }{caption}{fig:kl2}{0.9}

            % Same as λ L chain translation
            This process proceeds identically to lambda light chain translation, involving RER docking.

          \item \textbf{L REGION SNIPPING} [Mature Chain] --- \textit{ER Lumen}

            \pgfig{
              \hbandreset{0.2} % starting x-position
              \draw[dashed] (0,0) node[left] {N} -- (2.8,0) node[right] {C};

              \draw[hband pink={0.8}{V}];
              \draw[hband yellow={0.3}{J}];
              \draw[hband blue={1.2}{C}];

            }{caption}{fig:kl3}{0.9}

            % Enzymes clip off the leader, leaving a chain with a V(D)J domain and a constant domain
            Proteolytic removal of the leader sequence within the ER yields the final mature protein, comprising the variable and constant regions.

        \end{enumerate} 

      \end{core}

      \begin{core}{H Chain Gene Expression}

        \begin{enumerate}[label=(\arabic*)]

          \item \textbf{REARRANGEMENT} [Rearranged Gene] --- \textit{Nucleus}

            % First, a D region joins with a J, clipping out all extra D and J regions in between, but leaving any downstream Js
            The initial recombination event joins one Diversity (D) segment to one J segment. All intervening D and J segments are excised from the genome, while those situated downstream remain intact.

            \pgfig{
              \hbandreset{0.2}
              \draw[dashed] (0,0) node[left] {5'} -- (11,0) node[right] {\dots};

              \draw[hband black={0.2}{L}];
              \hbandshift{0.1}
              \draw[hband pink={0.8}{V}];
              \hbandshift{1.0}

              \draw[hband black={0.2}{L}];
              \hbandshift{0.1}
              \draw[hband pink={0.8}{V}];
              \hbandshift{0.6}

              \draw[hband green={0.2}{D}];
              \draw[hband yellow={0.3}{J}];
              \hbandshift{0.8}

              \draw[hband yellow={0.3}{J}];
              \hbandshift{0.8}

              \draw[hband yellow={0.3}{J}];
              \hbandshift{0.8}

              \draw[hband blue={1.2}{C}];
              \hbandshift{0.5}

              \draw[hband blue={1.2}{C}];
              \hbandshift{0.5}

            }{caption}{fig:lchain1}{0.85}

            % Second, a VL join with a D, removing all the downstream VLs and upstream Ds in between
            Following the D-J joining, a second recombination event splices a V segment to the newly formed D-J complex. This physically deletes all unselected V segments located downstream of the chosen V segment, and any D segments located upstream of the chosen D segment.

            \pgfig{
              \hbandreset{0.2}
              \draw[dashed] (0,0) node[left] {5'} -- (10.5,0) node[right] {\dots};

              \draw[hband black={0.2}{L}];
              \hbandshift{0.1}
              \draw[hband pink={0.8}{V}];
              \hbandshift{1.0}

              \draw[hband black={0.2}{L}];
              \hbandshift{0.1}
              \draw[hband pink={0.8}{V}];
              \draw[hband green={0.2}{D}];
              \draw[hband yellow={0.3}{J}];
              \hbandshift{0.8}

              \draw[hband yellow={0.3}{J}];
              \hbandshift{0.8}

              \draw[hband yellow={0.3}{J}];
              \hbandshift{0.8}

              \draw[hband blue={1.2}{C}];
              \hbandshift{0.5}

              \draw[hband blue={1.2}{C}];
              \hbandshift{0.5}

            }{caption}{fig:lchain2}{0.85}

            \pgfig{
              \hbandreset{0.2}
              \draw[dashed] (0,0) node[left] {5'} -- (14.5,0) node[right] {\dots};

              \draw[hband black={0.2}{L}];
              \hbandshift{0.1}
              \draw[hband pink={0.8}{V}];
              \hbandshift{5}
              \draw[hband green={0.2}{D}];
              \hbandshift{5}
              \draw[hband yellow={0.3}{J}];
              \hbandshift{0.8}

              \draw[hband blue={1.2}{C}];
              \hbandshift{0.5}

              \draw[fill=imgray, draw=none] (1.3,-0.3) rectangle (2.0,0.3);
              \draw[fill=impurple, draw=none] (2.0,-0.3) rectangle (2.7,0.3);
              \draw[fill=imorange, draw=none] (3.3,-0.3) -- (2.7,-0.3) -- (2.7,0.3) -- cycle;
              \draw (2.7,0.3) -- (1.3,0.3) -- (1.3,-0.3) -- (3.3,-0.3) -- cycle; 

              \draw[fill=imlightpurple, draw=none] (4.9,-0.3) rectangle (5.6,0.3);
              \draw[fill=imgray, draw=none] (5.6,-0.3) rectangle (6.3,0.3);
              \draw[fill=imorange, draw=none] (4.3,0.3) -- (4.9,0.3) -- (4.9,-0.3) -- cycle;
              \draw (4.3,0.3) -- (6.3,0.3) -- (6.3,-0.3) -- (4.9,-0.3) -- cycle; 

              \draw[fill=imgray, draw=none] (6.5,-0.3) rectangle (7.2,0.3);
              \draw[fill=imlightpurple, draw=none] (7.2,-0.3) rectangle (7.9,0.3);
              \draw[fill=imorange, draw=none] (7.9,-0.3) -- (7.9,0.3) -- (8.5,0.3) -- cycle;
              \draw (8.5,0.3) -- (6.5,0.3) -- (6.5,-0.3) -- (7.9,-0.3) -- cycle; 

              \draw[fill=impurple, draw=none] (10.1,-0.3) rectangle (10.8,0.3);
              \draw[fill=imgray, draw=none] (10.8,-0.3) rectangle (11.5,0.3);
              \draw[fill=imorange, draw=none] (9.5,-0.3) -- (10.1,-0.3) -- (10.1,0.3) -- cycle;
              \draw (10.1,0.3) -- (11.5,0.3) -- (11.5,-0.3) -- (9.5,-0.3) -- cycle;

            }{caption}{fig:rss3}{0.85}

          \item \textbf{TRANSCRIPTION} [Primary Transcript] --- \textit{Nucleus}

            \pgfig{
              \hbandreset{0.2}
              \draw[dashed] (0,0) node[left] {5'} -- (10.5,0) node[right] {\dots};

              \draw[hband black={0.2}{L}];
              \hbandshift{0.1}
              \draw[hband pink={0.8}{V}];
              \hbandshift{1.0}

              \draw[hband black={0.2}{L}];
              \hbandshift{0.1}
              \draw[hband pink={0.8}{V}];
              \draw[hband green={0.2}{D}];
              \draw[hband yellow={0.3}{J}];
              \hbandshift{0.8}

              \draw[hband yellow={0.3}{J}];
              \hbandshift{0.8}

              \draw[hband yellow={0.3}{J}];
              \hbandshift{0.8}

              \draw[hband blue={1.2}{\ce{C_μ}}];
              \hbandshift{0.5}

              \draw[hband blue={1.2}{\ce{C_δ}}];
              \hbandshift{0.5}

            }{caption}{fig:lchain3}{0.85}

            % RNA polymerase transcribes from the LVDJ region, through any remaining Js and introns, and copies through the constant regions and stops at the first stop signal at the end of Cδ (after passing through Cμ)
            The primary transcript synthesized spans the fully rearranged V-D-J sequence, encompassing all downstream introns and remaining J segments. Importantly, transcription continues unabated through both the Cμ and Cδ constant region coding domains, terminating only after the poly-A site situated downstream of Cδ.

          \item \textbf{PROCESSING} [Processed Message] --- \textit{Nucleus}

            \pgfig{
              \hbandreset{0.2}
              \draw[dashed] (0,0) node[left] {5'} -- (3.5,0) node[right] {AA(A\ce{_n})AA 3'};

              \draw[hband black={0.2}{L}];
              \draw[hband pink={0.8}{V}];
              \draw[hband green={0.2}{D}];
              \draw[hband yellow={0.3}{J}];
              \draw[hband blue={1.2}{\ce{C_μ}}];

              \hbandreset{0.2}
              \begin{scope}[shift={(7,0)}]
                \draw[dashed] (0,0) node[left] {5'} -- (3.5,0) node[right] {AA(A\ce{_n})AA 3'};

                \draw[hband black={0.2}{L}];
                \draw[hband pink={0.8}{V}];
                \draw[hband green={0.2}{D}];
                \draw[hband yellow={0.3}{J}];
                \draw[hband blue={1.2}{\ce{C_δ}}];
              \end{scope}

            }{caption}{fig:lchain4}{0.85}

            % The transcript undergoes alternative message splicing to include either the μ or the δ constant exon, but not both
            The elongated primary transcript is subject to alternative RNA splicing, a critical regulatory step determining whether the mature mRNA will ultimately incorporate the exons for the μ heavy chain (forming IgM) or the δ heavy chain (forming IgD), but not both simultaneously on the same transcript.

            % Any introns and extra Js get clipped out
            Concurrently, all intronic sequences and superfluous unselected J segments are spliced out.

          \item \textbf{TRANSLATION} [Nascent Peptide] --- \textit{Into the RER}

            \pgfig{
              \hbandreset{0.2}
              \draw[dashed] (0,0) node[left] {N} -- (3.5,0) node[right] {C};

              \draw[hband black={0.2}{L}];
              \draw[hband pink={0.8}{V}];
              \draw[hband green={0.2}{D}];
              \draw[hband yellow={0.3}{J}];
              \draw[hband blue={1.2}{\ce{C_μ}}];

              \hbandreset{0.2}
              \begin{scope}[shift={(5,0)}]
                \draw[dashed] (0,0) node[left] {N} -- (3.5,0) node[right] {C};

                \draw[hband black={0.2}{L}];
                \draw[hband pink={0.8}{V}];
                \draw[hband green={0.2}{D}];
                \draw[hband yellow={0.3}{J}];
                \draw[hband blue={1.2}{\ce{C_δ}}];
              \end{scope}

            }{caption}{fig:lchain5}{0.85}

            % Same as L chain translation
            Follows the same ER-docking and luminal translocation process described for light chains.

          \item \textbf{L REGION SNIPPING} [Mature Chain] --- \textit{ER Lumen}

            \pgfig{
              \hbandreset{0.2}
              \draw[dashed] (0,0) node[left] {N} -- (3.5,0) node[right] {C};

              \draw[hband pink={0.8}{V}];
              \draw[hband green={0.2}{D}];
              \draw[hband yellow={0.3}{J}];
              \draw[hband blue={1.2}{\ce{C_μ}}];

              \hbandreset{0.2}
              \begin{scope}[shift={(5,0)}]
                \draw[dashed] (0,0) node[left] {N} -- (3.5,0) node[right] {C};

                \draw[hband pink={0.8}{V}];
                \draw[hband green={0.2}{D}];
                \draw[hband yellow={0.3}{J}];
                \draw[hband blue={1.2}{\ce{C_δ}}];
              \end{scope}

            }{caption}{fig:lchain6}{0.85}

            % Enzymes clip off the leader, leaving a chain with a V(D)J domain and a constant domain
            Proteolytic removal of the leader sequence yields the mature heavy chain polypeptide, featuring the highly variable V-D-J domain and the selected constant isotype domain.

        \end{enumerate} 

      \end{core}

      \begin{core}{RSS Joining Mechanism}

        \begin{enumerate}

          \item \textbf{Juxtaposition} 

            \begin{itemize}

              \item \textit{Locating \& Cleaving}

                % The genes to be arranged move to specific locations in the nucleus and open up, detaching from the nucleosomes (organized DNA units)
                Target loci destined for rearrangement physically relocate to transcriptionally active zones within the nucleus. The chromatin structure relaxes, displacing nucleosomes to grant enzymatic access to the naked DNA.

                % The processing complex grabs one of the RSS signals—one-turn juxtaposes with two-turn—and one strand of the DNA in 5’ heptamer gets cleaved
                The RAG-1/RAG-2 recombinase complex aligns specific Recombination Signal Sequences (RSS), strictly adhering to the "12/23 rule," where a sequence featuring a one-turn spacer must pair with a sequence featuring a two-turn spacer. Upon alignment, an endonucleolytic nick is introduced precisely at the 5' border of the conserved heptamer sequence.

                <FIGURE>

              \item \textit{Hairpin Formation}

                % This loose 5’ end heptamer attacks the opposite side of the uncut strand of the same DNA molecule, producing a hairpin loop at the downstream end of the V and the upstream end of the J (light chains) or D (heavy chains)
                The newly liberated 3' hydroxyl group executes a nucleophilic attack on the complementary phosphodiester bond of the intact strand. This transesterification reaction creates a covalently sealed hairpin loop at the coding ends of the participating V and J segments (in light chains), or the V, D, and J segments (in heavy chains).

                % ***The spacer sequences, along with the AT region signal end in a flush cut with a 5’ phosphorylated end
                Conversely, the excised signal ends—containing the heptamer, spacer, and nonamer sequences—are left as a blunt, flush cut featuring a free 5' phosphate group.

                <FIGURE>

            \end{itemize}

          \item \textbf{Ligation}

            \begin{itemize}

              \item \textit{Initial Cleavage \& Ligation}

                <FIGURE>

                % Enzymes clip the hairpin loop at random locations, ligating two of the single strands together and leaving a region of unmatched single stranded regions
                The Artemis nuclease complex resolves the coding hairpins through an asymmetric, random single-stranded cleavage. This imprecise nicking generates varying lengths of unmatched single-stranded DNA overhangs at the joining interface.

              \item \textit{P-Nucleotide Addition}

                <FIGURE>

                % Nucleotides are added to fill in the unmatched single stranded regions, matching the ends generated by the cut
                Cellular repair polymerases fill in these asymmetrical overhangs, synthesizing short palindromic (P) nucleotide sequences that conform perfectly to the templates created by the offset hairpin resolution.

              \item \textit{N-Nucleotide Addition (Heavy Chains Only)}

                <FIGURE>

                % Up to 15 nucleotides can be added at random in the junction (further adding to the junction’s variability)
                In developing heavy chains, the enzyme Terminal Deoxynucleotidyl Transferase (TdT) can independently append up to 15 non-templated (N) nucleotides randomly into the junctional gap prior to final ligation, profoundly expanding the sequence diversity of the resulting antigen receptor.

            \end{itemize}

          \item \textbf{Result}

            % Thus at the end, there is repair and ligation of the coding joint (coding regions + heptamers), and release and digestion of the signal sequences (spacer + AT-nonamer + portion of heptamer; which are retained if there’s an inversional joining)
            The culmination of this process yields a fully repaired and ligated coding joint, forming the functional antigen-binding domain. Meanwhile, the excised signal sequence fragment—typically forming a closed episomal circle consisting of the joined heptamer-spacer-nonamer elements—is generally degraded, unless specific structural topologies dictated an inversional joining event.

            % Somatic Hypermutation—at a later stage of development, the parts of the genes coding for the hypervariable region can mutate, introducing further variation
            Subsequent to baseline V(D)J recombination, mature B cells undergo Somatic Hypermutation. During active immune responses, point mutations are intentionally introduced into the hypervariable gene loci, providing the substrate for affinity maturation.

            % ***This occurs later in the process of antibody stimulation, maturation, and selection outside the bone marrow
            It is essential to delineate that somatic hypermutation occurs exclusively in peripheral secondary lymphoid tissues following antigen stimulation, well after the initial B cell maturation phase in the bone marrow has concluded.

        \end{enumerate}

      \end{core}

  \subsection{Other}

    \begin{core}{\textit{Edelman \& Porter}}

      In foundational experiments conducted during the 1960s, researchers utilizing serum protein electrophoresis fractionated human plasma into five distinct molecular cohorts. One specific migratory band, designated the "\textit{gamma}" globulin fraction, was identified as housing the vast majority of the plasma's protective immunological activity. Subsequent purification of this fraction \textbf{led to the biochemical isolation and characterization of antibody molecules}.

      \begin{enumerate}

        \item \textbf{Gerald Edelman} (\textit{Mercaptoethanol})

          Edelman employed \textbf{mercaptoethanol} to \textbf{reduce the disulfide bonds} that keep the light and heavy chains interconnected with one another, thereby enabling each chain to be isolated and studied individually.

          \myfig{0.7}{immunology_figures/edelman}{Antibody cleavage via mercaptoethanol; \textit{cleavage occurs at disulfide bonds, reducing the antibody to individual light and heavy chains}}{fig:edelman}

        \item \textbf{Rodney Porter} (\textit{Pepsin \& Papain})

          Porter demonstrated that \textbf{pepsin} cleaves preferentially at the \textbf{bottom of the hinge} between the arms and the stem, thereby disconnecting the two individual heavy chain Fc sections from the top half of the molecule and yielding a single \textbf{Fab2 fragment}.

          \myfig{0.65}{immunology_figures/porter1}{Antibody cleavage via pepsin; \textit{cleavage occurs at hinge, separating two heavy chain fragments and one Fab2 molecule}}{fig:porter1}

          By contrast, he found that \textbf{papain} cleaves preferentially at the \textbf{top of the hinge} between the arms and the stem, disconnecting the Fc sections---which remain joined to one another by disulfide bonds---from two separate \textit{Fab} fragments.

          \myfig{0.6}{immunology_figures/porter2}{Antibody cleavage via papain; \textit{cleavage occurs at upper hinge, resulting in connected lower heavy chain fragments and two Fab fragments}}{fig:porter2}

          Notably, structurally \textbf{divalent Fab2 fragments} (\textit{from pepsin treatement}) remain capable of precipitating or cross-linking antigens to form extended immune complexes, whereas \textbf{monomeric Fab fragments} (\textit{from papain treatment}) lack this cross-linking ability.

      \end{enumerate}

    \end{core}

    \begin{core}{Haptens (\textit{Small Antigens})}

      \textbf{Haptens} are \textbf{low-molecular-weight} chemical entities that can be recognized by antibodies but are \textbf{inherently non-immunogenic}; their minute size prevents them from independently eliciting an adaptive immune response.

      To overcome this limitation for experimental or clinical antibody production, haptens must be covalently conjugated to a \textbf{high-molecular-weight} carrier protein, such as \textbf{bovine serum albumin (BSA)}, to provide the requisite steric bulk and T cell epitopes needed to induce active immunity.

      <FIGURE>

    \end{core}

    \begin{core}{Monoclonal Antibodies (mAb)}

      % Having an endless production stream of an antibody with a desired CDR for a certain antigen epitope has broad applications in the pharmaceutical industry—particularly in the field of oncology.
      The ability to engineer continuous, uniform production of antibodies featuring customized, highly specific CDR loops targeting defined antigen epitopes holds profound clinical and pharmacological value, especially concerning targeted oncological therapeutics.

      % Definition: Monoclonal antibodies are fractions of antibodies with an identical defined specificity (CDR region) for a particular antigen
      Definition: Monoclonal antibodies represent a homogenous population of immunoglobulins synthesized by a single clonal lineage of B cells, resulting in absolute uniformity of paratope structure and antigen specificity.

      % Hybridomas are a species of these B cells that are immortal and will endlessly crank out a pure stream of antibodies for you. Here’s how you make it.
      This production is achieved through the generation of hybridomas—engineered, immortalized cell lines formed by fusing specific primary B cells with myeloma cells, thereby ensuring an indefinite, robust yield of the desired monoclonal antibody construct. The general synthesis methodology is detailed below.

      <FIGURE>

    \end{core}

\section{B Cells}

  % \subsection{B Cell Receptors}

  %   \begin{core}{B Cell Receptors}

  %     \sbs{0.45}{
  %       A \textbf{B cell receptor (BCR)} is an immunoglobulin molecule anchored to the cellular membrane and oriented outward, which recognizes foreign antigen when two neighboring receptors bind it and undergo \textbf{crosslinking}. Because the BCR itself has only a very short cytoplasmic tail incapable of independent signaling, it relies entirely on the associated heterodimer of Igα/Igβ---currently designated CD79a and CD79b, respectively---to relay the crosslinking event into an intracellular signaling cascade, ultimately triggering downstream pathways.

  %       Prior to activation, naive B cells exclusively co-express IgM and IgD as their surface receptors, both of which share the same antigen-binding specificity despite differing in their heavy chain constant regions. Following activation and differentiation, memory B cells instead display receptors belonging to other heavy chain classes (\textit{e.g., IgG, IgA, or IgE}), a transition driven by exposure to antigen and helper T cell signals.
  %     }{
  %       \myfig{0.98}{immunology_figures/bcr}{B cell receptor}{fig:bcr}
  %     }

  %   \end{core}

  \subsection{Naive B Cell Development}

    \begin{core}{Naive B Cell Development}

      \begin{enumerate}[label=(\arabic*)]

        \item \textbf{Co-Receptor Development}

          <FIGURE>

          % Exposing CLPs to B-Cell Specific Activator Protein (BSAP) triggers the formation of BCR co-receptors (Igα/Igβ) on the cell surface, pushing the cell down the B cell pathway.
          The exposure of Common Lymphoid Progenitors (CLPs) to B-Cell Specific Activator Protein (BSAP) induces the initial surface expression of essential BCR signaling co-receptors, specifically Ig$\alpha$ and Ig$\beta$, firmly committing the precursor to the B lymphocyte developmental trajectory.

          % This cell is called a pro B cell.
          At this stage of commitment, the developing cell is classified as a progenitor B (pro-B) cell.

          % BSAP is present in ALL members of the B cell lineage except mature plasma cells, which are done differentiating
          BSAP expression is maintained constantly across the entire B cell lineage continuum, ceasing only upon terminal differentiation into antibody-secreting plasma cells.

        \item \textbf{Heavy Chain}

          <FIGURE>

          % The pro B cell rearranges a H chain gene (max two tries or apoptosis)
          The pro-B cell initiates V(D)J rearrangement at the heavy chain locus. The cell is afforded a maximum of two opportunities (one per allele) to generate a productive rearrangement before default apoptosis occurs.

          % After successful rearrangement and formation, the H chain peptide picks up a surrogate light chains, becoming a pre B cell 
          Upon translating a viable heavy chain polypeptide, the nascent heavy chain associates with a pre-existing surrogate light chain complex, advancing the cell to the precursor B (pre-B) cell stage.

          % The attachment of this fake L chain signals allelic exclusion, serves as a temporary filler for the soon-to-be produced L chain, and initiates light chain gene rearrangement 
          The functional assembly of this surrogate receptor complex delivers a crucial intracellular signal that enforces allelic exclusion at the heavy chain locus, maintains structural stability, and signals the nucleus to begin rearranging the endogenous light chain genes.

        \item \textbf{Light Chain}

          <FIGURE>

          % The pre B cell rearranges a L chain gene (max four tries or apoptosis)
          The pre-B cell subsequently attempts V-J rearrangement of the light chain loci, possessing up to four total opportunities (spanning both the $\kappa$ and $\lambda$ alleles) before risking apoptosis.

          % If there is successful rearrangement, the receptors then will have displayed an IgM BCR with a determined antigenic specificity (CDR)
          A productive rearrangement yields a functional light chain that pairs with the heavy chain, resulting in the surface expression of a complete IgM BCR featuring a fully defined complementarity-determining region (CDR).

          % This cell is called an immature B cell.
          The cell, now bearing its unique surface receptor, transitions into the immature B cell designation.

        \item \textbf{Negative Selection}

          % The immature B cell will undergo a process to see whether the receptors will attach to self-antigens (negative selection)
          Within the bone marrow microenvironment, immature B cells undergo negative selection to rigorously test their newly formed BCRs for potentially harmful autoreactivity against local self-antigens.

            % Only about 10\% of the B cells will make it through this process
            Approximately 10\% of the initial developing B cell pool successfully navigates this strict tolerance checkpoint.

            % Elimination is signaled by the crosslinking of IgM, leading to apoptosis
            Autoreactive cells face elimination; binding of widespread multivalent self-antigens heavily crosslinks the surface IgM, triggering rapid apoptotic cell death pathways.

            % HOWEVER, some cells can sometimes get a second chance, reactivating their RAG (Recombination Activating Gene) enzymes, and try rearranging again
            However, under specific conditions, an autoreactive B cell may avoid immediate destruction by reactivating Recombination Activating Gene (RAG) transcription, initiating a process termed receptor editing to generate a new, non-autoreactive light chain specificity.

          % Successful cells do not recognize self-antigens
          Cells that successfully bypass negative selection exhibit functional central tolerance, lacking recognition for endogenous self-antigens.

        \item \textbf{IgD Formation}

          <FIGURE>

          % The cell now has the option to splice the BCR message so that it can synthesize the IgD version of its BCRs
          Surviving cells employ alternative mRNA splicing to concurrently transcribe and translate the IgD isotype alongside their existing IgM receptors.

          % This is now a mature B cell and can leave the bone marrow
          The dual expression of surface IgM and IgD signifies a fully mature, immunocompetent, yet naive B cell, which is now authorized to exit the bone marrow and enter systemic circulation.

        \item \textbf{During Circulation}

          % The cells are released into the plasma and head to the peripheral lymphoid organs: lymph nodes, spleen, and mucosa
          The mature B cells migrate through the circulatory and lymphatic systems to seed peripheral secondary lymphoid organs, including the lymph nodes, spleen, and various mucosal tissues.

          % Once stimulated to produce antibodies, the cell makes primarily μ peptides, but without the membrane spanning region (thus they will make the soluble IgM pentamer)
          Upon encountering specific antigen and receiving activating signals, the resulting plasma cells drastically upregulate transcription of the $\mu$ heavy chain, altering RNA processing to exclude the transmembrane coding domains, leading to massive secretion of the soluble, multimeric IgM pentamer.

          % The cell rarely makes δ peptides without the membrane spanning regions.
          Soluble secretion of the $\delta$ isotype (IgD) is physiologically negligible, as the cell rarely produces $\delta$ transcripts lacking the membrane anchor sequences.

      \end{enumerate}

    \end{core}

  \subsection{B Cell Activation and Proliferation}

    \subsubsection{Thymus Independent and Dependent B Cell Activation}

      \begin{core}{Thymus Independent B Cell Activation}

        % There exists few antigens (TI antigens) that can prompt B-cell development independent of TH cell co-stimulus. These can also simultaneously activate TLRs
        Certain distinct classes of antigens, termed Thymus-Independent (TI) antigens, possess the unique capability to induce B cell activation and proliferation in the complete absence of concurrent T helper (TH) cell stimulation. These structures frequently co-engage innate Toll-Like Receptors (TLRs).

        % TI activation [1] does not induce class switching (you make IgM) and [2] does not produce memory cells; for either you need TH cells
        TI-mediated activation is characterized by two fundamental limitations: it generally fails to induce functional immunoglobulin class switching (restricting output predominantly to IgM), and it does not yield a sustained population of memory B cells; robust induction of either process strictly requires classical TH cell involvement.

        \begin{enumerate}

          \item \textbf{Type 1} 

            <FIGURE>

            % Lipopolysaccharides such as those found in the cell wall of gram negative bacteria 
            TI-1 antigens typically consist of potent mitogenic molecules, most notably lipopolysaccharide (LPS), derived from the outer leaflet of Gram-negative bacterial cell walls.

            % These molecules also activate TLR4 to provide the second permissive signal
            At significant concentrations, these molecules bind the BCR while simultaneously acting as potent ligands for adjacent TLR4 receptors, delivering the necessary, self-contained secondary permissive signal for cellular activation.

          \item \textbf{Type 2}

            <FIGURE>

            % Repetitive polymeric proteins (such as bacterial flagellin) cross-link and clusters Ig receptors
            TI-2 antigens are characterized by highly repetitive, polymeric structural motifs (such as extensive bacterial capsular polysaccharides or flagellin polymers) that induce massive, sustained cross-linking and clustering of surface BCRs.

            % The signal does not come from a discrete source, but inflammatory cytokines (IL-2, IL3, and IFNɣ) and APCs
            While lacking direct TH cell contact, successful TI-2 activation often relies upon a generalized, supportive microenvironment heavily supplemented by inflammatory cytokines (e.g., IL-2, IL-3, and IFN-$\gamma$) secreted by nearby innate immune populations and specialized APCs.

        \end{enumerate}

    \end{core}

    \begin{core}{Thymus Dependent B Cell Activation}

      % IN MOST CASES, antigen cross-links a B cell’s Ig receptors, and sets off a signal, making the B cell seek out a TH cell.
      Under canonical circumstances governing the majority of adaptive responses, specific antigen binding heavily cross-links the B cell’s surface immunoglobulin. This initial binding event triggers an intracellular cascade that structurally primes the B cell and directs it to actively solicit requisite co-stimulation from an antigen-specific TH cell.

      <FIGURE>

      \begin{enumerate}[label=(\arabic*)]

        \item % Antigen cross links BCR Ig, bringing together cytosolic Igα/Igβ domains and activating the ITAMs
        Extracellular antigen physically cross-links multiple BCR complexes. This physical clustering forces the spatial aggregation of the associated intracellular Ig$\alpha$/Ig$\beta$ signaling domains, leading to the rapid structural activation of their Immunoreceptor Tyrosine-based Activation Motifs (ITAMs).

        \item % Src-like kinases attach to ITAMs, adding phosphate from ATP
        Resident Src-family kinases bind these aggregated ITAMs, utilizing ATP hydrolysis to systematically phosphorylate critical tyrosine residues on the motifs.

        \item % Syk kinase docks and sets off a complex interlocking activation cascade
        The newly phosphorylated ITAMs serve as high-affinity docking sites for Syk kinase. The recruitment and activation of Syk subsequently initiates a vast, highly interconnected intracellular signaling web:

          \begin{enumerate}

            \item % DAG-PIP: An up-regulated hydrolase cleaves a membrane phospholipid to diacyl glycerol (DAG) and phosphoinositide (PIP)—trigger many responses 
            DAG/IP3 Generation: Phospholipase C-$\gamma$ (PLC-$\gamma$) is activated, cleaving the membrane phospholipid PIP2 into diacylglycerol (DAG) and inositol trisphosphate (IP3), acting as pivotal second messengers.

              % → release of Ca2+ from the ER → activates calmodulin → NFAT (active transcription factors and upregulates genes) activation
              IP3 diffuses to trigger the rapid release of stored Ca2+ from the ER lumen into the cytosol. This calcium influx heavily binds calmodulin, which subsequently activates the phosphatase calcineurin to dephosphorylate and activate the transcription factor NFAT.

              % → removal of iκB (inhibitor) → up-regulation of inflammatory transcription factor NF-κB
              Concurrently, DAG activation propagates signals that result in the targeted degradation of the I$\kappa$B inhibitory complex, thereby unleashing the potent pro-inflammatory transcription factor NF-$\kappa$B to enter the nucleus.

            \item % RAS-RAF-MEK-ERK kinases: This pathway leads to many activated ERKs, causing them to activate, bind DNA and turn on genes for B cell division and differentiation
            MAPK/ERK Kinase Cascade: Through intermediary adapter proteins, the classical RAS-RAF-MEK-ERK phosphorylation cascade is initiated. This culminates in the nuclear translocation of activated ERK proteins, which directly upregulate essential target genes controlling robust cellular proliferation and terminal differentiation.

            \item % PIP2-PIP3: activation of the BCR causes the PIP2 to add a phosphate and turn into PIP3, blocking signals leading to apoptosis
            PI3K Activation: The activated BCR complex recruits PI3K, which phosphorylates residual membrane PIP2 to yield PIP3. The accumulation of PIP3 provides a powerful overarching survival signal, effectively blocking endogenous apoptotic pathways during the stress of rapid clonal expansion.

          \end{enumerate}

        \item % The cells begin to divide and secrete antibodies
        Driven by these integrated transcription factors, the fully activated B cells commence exponential mitotic division and begin differentiating into functional, antibody-secreting plasma cells.

        \item % REGULATION—As antibody-antigen complexes build up, they bind to the Fc receptors on the B cells (CD22), which provides brakes on the system (FIGURE)
        Feedback Regulation: As the immune response progresses, the accumulating high-titer antigen-antibody immune complexes begin to bind inhibitory surface Fc$\gamma$ receptors, specifically CD32 (Fc$\gamma$RIIB) which houses an ITIM domain, providing essential negative feedback to safely curtail the response. (Note: CD22 is an independent inhibitory coreceptor recognizing sialic acid, playing a similar dampening role).

          \begin{enumerate}

            \item % Developing tolerance for foreign antigens refers to the immune system's ability to recognize and accept the presence of specific antigens without mounting an immune response against them—do not want this against pathogenic infections
            While strictly necessary for controlling inflammation, suppressing active signaling presents a risk of inducing peripheral tolerance toward foreign antigens—a state where the immune system inappropriately acknowledges a foreign entity without mounting a defense. This is heavily counterproductive during active pathogenic infections.

            \item % Treg cells further suppresses immune responses to your own proteins and to those of benign commensal bacteria and fungi
            To carefully manage this balance, specialized Regulatory T (Treg) cells act primarily to enforce systemic tolerance specifically directed against autologous host proteins, as well as maintaining non-reactivity toward benign, commensal microbiome flora.

            \item % IF YOU INTRODUCE FOREIGN ANTIBODIES to an antigen, this will tie up the antigen and prevent it from promoting an immune response
            Clinically, if exogenous antibodies targeting a specific antigen are passively administered into a host system, they will effectively sequester the circulating antigen load, largely preventing it from natively binding host BCRs and thus actively suppressing the generation of a de novo immune response.

          \end{enumerate}

      \end{enumerate}

    \end{core}

    \begin{core}{T$_\text{H}$ Cell Synpase}

      % The BCR does not signal nor divide effectively without contact with a TH cell—the B cell can stimulate a TH cell at concentrations 100 to 10,000 times lower than those necessary for a macrophage or dendritic cell
      Despite robust BCR engagement, sustained proliferation heavily relies upon physical contact with an activated TH cell. Notably, the high affinity of the BCR allows the B cell to capture and present antigen to TH cells at concentrations roughly 100 to 10,000 times lower than what is required for standard presentation by macrophages or dendritic cells.

      \begin{enumerate}[label=(\arabic*)]

        \item % When B cells bind antigen, they bring some inside and hydrolyze it
        BCR engagement triggers receptor-mediated endocytosis, allowing the B cell to internalize the bound antigen and subsequently hydrolyze it into small peptide fragments via the endosomal pathway.

        \item % The processed antigen is expressed through the MHC II on the B cell surface
        These processed antigenic peptides are then loaded onto nascent MHC class II molecules and shuttled back to the external cell surface for presentation.

        \item % The [B cell] MHC II-antigen complex binds to the [TH cell] TCR with the help of CD4, forming a immune synapse
        An activated TH cell specific for that exact peptide epitope binds the MHC II complex via its TCR, heavily stabilized by the CD4 co-receptor, forming a highly organized cellular interface termed the immunological synapse.

        \item % This causes the TH cell to produce CD40L (or CD154) which binds to the surface CD40R on the B cell
        This direct recognition event induces the TH cell to rapidly upregulate surface expression of CD40 Ligand (CD40L/CD154), which securely engages the constitutively expressed CD40 receptor physically present on the B cell surface.

        \item % This leads to the production of the second signal, activating the interlocking up-regulatory pathways of the B cell and they begin proliferating and differentiating
        The CD40-CD40L interaction delivers the critical, definitive second signal to the B cell. This massively amplifies the internal signaling cascades initiated by the BCR, fully committing the B cell to exponential proliferation and terminal effector differentiation.

      \end{enumerate}

      % In the reciprocal part of the conversation between B and T cell, B7 plays a role in the up-regulation of the TH cell
      In this reciprocal exchange, the B cell concurrently upregulates surface B7 molecules (CD80/CD86), which engage CD28 on the TH cell to continually support and reinforce the TH cell’s activated state.

    \end{core}

  \subsubsection{Antibody Development}

    \begin{core}{Lymph Node Follicle Development}

      <FIGURE>

      % Free and antibody-bound antigen in the plasma is likely to be picked up by APCsj
      Circulating free antigen and small immune complexes present in the plasma or lymph are predominantly captured by local antigen-presenting cells (APCs).

      % APCs initially activate the TH cells in the paracortex which then goes on to synapse with B cells, producing a clonal cluster at the boundary between the paracortex and cortex
      These APCs traffic to the lymph node to initially activate TH cells within the paracortical region. These primed TH cells then migrate toward the B cell zones, forming functional immunological synapses with antigen-bearing B cells. This interaction generates a proliferating cellular cluster localized precisely at the boundary delineating the paracortex and the outer cortex.

      % This eventually morphs into a secondary follicle (from a primary follicle) with a light germinal center and dark zone
      This active cluster rapidly expands and reorganizes, driving the transformation of a resting primary follicle into a highly active secondary follicle, distinctly characterized by a well-defined light germinal center and a dense dark zone.

        % Germinal Center: Where B, TH , and follicular DCs interact for antibody selection 
        Germinal Center: The primary structural environment facilitating the intense interactions among B cells, specialized TH cells, and FDCs required for rigorous antibody selection.

        % Dark Zone: Sites of mutation and excess reproduction of B cells
        Dark Zone: The anatomical location dedicated to extreme B cell proliferation accompanied by targeted somatic hypermutation.

      % REMINDER: FDCs capture antigen-antibody complexes in beaded structures (iccosomes) and present them to B cells
      NOTE: Within this environment, FDCs trap and stabilize intact antigen-antibody immune complexes within distinct beaded cellular structures termed iccosomes, presenting them efficiently for selection by competing B cell clones.

    \end{core}

    \begin{core}{Affinity Maturation}

      <FIGURE>

      % Centroblasts (activated, enlarged B cells) move to the dark zone where they undergo [1] rapid growth and division and [2] somatic hypermutation (CDR loop mutations)—during this phase, they stop displaying their membrane Ig receptors
      Centroblasts—highly activated, morphologically enlarged B cells—migrate into the dark zone where they experience intense mitotic proliferation. Concurrently, they undergo somatic hypermutation, introducing targeted nucleotide substitutions strictly within the DNA sequences coding for the CDR loops. Throughout this proliferative stage, centroblasts temporarily cease the surface expression of their BCRs.

      % Centrocytes (centroblasts that stop dividing) move to the germinal center or light zone to undergo selection where B cells with more effective receptors will survive and multiply at a greater rate by connecting with iccosomes on the FDCs and receiving CD40 signals from TH cells (ACCELERATED NATURAL SELECTION)
      Upon exiting the cell cycle, these cells mature into centrocytes and migrate toward the light zone for a rigorous selection process. Centrocytes must actively compete to bind the limited antigen presented on FDC iccosomes. Cells possessing mutated receptors with improved affinity successfully bind the antigen, enabling them to present peptide to TH cells, thereby securing the vital CD40 survival signals. This process functions essentially as accelerated, localized natural selection.

        % Centrocytes that don’t contact an antigen die by apoptosis and are recycled by macrophages
        Centrocytes bearing low-affinity or non-functional mutations fail to secure sufficient antigen or TH cell help; consequently, they undergo rapid apoptosis and are cleared locally by tingible body macrophages.

      % The centrocytes leave the germinal center as plasma cells, lose their surface antibody again, and begin secreting antibodies at the site of infection
      The highly selected, high-affinity centrocytes finally exit the germinal center microenvironment, terminally differentiating into long-lived plasma cells. These cells permanently downregulate surface receptor expression in favor of massive, sustained secretion of high-affinity antibodies targeted against the infection.

    \end{core}

    \begin{core}{Class Switching}

      % TH cells synapse with the B cells and cytokine signal the specific switch regions (recombination sites designated by 2-3 KB sequences of DNA)
      TH cells maintain synaptic contact with differentiating B cells, secreting tailored cytokine profiles that chemically direct class-switch recombination. This mechanism physically targets specific "switch regions"—conserved sequences spanning 2 to 3 kilobases positioned immediately upstream of each respective constant gene segment.

      % ***Before switching, the cell would be expressing the µ and δ constant regions
      Prior to any switching event, the mature B cell exclusively transcribes the proximal $\mu$ and $\delta$ constant regions, forming its initial IgM/IgD receptors.

      % Switching involves removing these and whatever other C gene region(s) stand(s) between the VDJ region and the C region to be expressed
      The switching mechanism involves the permanent excisional recombination of DNA, physically looping out and discarding the previously utilized constant domains alongside any intervening sequences positioned between the rearranged VDJ domain and the newly targeted downstream constant region.

      % Any C region removed upstream after class switching is unretrievable, thus removing the last Cα region will not allow a cell to produce any antibody at all
      Because this deletion occurs at the DNA level, any excised upstream constant regions are irretrievably lost from that allele's genome. Therefore, progressive class switching is unidirectional down the heavy chain locus.

      % Takes place after the system has been producing antibodies for a week or more
      This large-scale genetic remodeling generally requires a sustained immune response, typically occurring functionally prominent after a week or more of persistent antigen challenge.

      % ***First antibodies produced for an infection are IgMs and you don’t start seeing IgGs (or other classes) until later in
      Consequently, the primary serological hallmark of a novel acute infection is a rapid spike in low-affinity IgM; the transition to high-affinity IgG (or other specialized isotypes) is a slower, distinctly subsequent event characterizing the later phases of the adaptive response.

      <FIGURE>
      
    \end{core}

  \subsubsection{Plasma vs Memory B Cells}

    \begin{core}{Plasma Cells}

      % The final differentiation to a plasma cell (from centrocytes) involve the switch (Blimp-1, XBP-1, IRF-4 from ChatGPT) that generates the splicing enzymes that do not add the membrane-spanning exons to the μ heavy chain message
      The terminal differentiation from an activated centrocyte into a highly specialized plasma cell is orchestrated by key transcription factors. This shifts the cellular machinery, altering RNA splicing enzymes to selectively exclude the membrane-anchoring exons from the heavy chain transcript, thereby converting receptor expression into large-scale soluble secretion.

      % Low levels of BSAP (differentiating related factor) tend to promote plasma cell formation
      This terminal transition is strongly correlated with, and facilitated by, the significant downregulation and suppression of BSAP (Pax-5) expression.

    \end{core}

    \begin{core}{Memory Cells}

      % Memory cells set aside from this process may resemble naive B cells, but they have undergone class switching, thus they express a larger variety of BCRs compared to the naive B cell (only IgD and IgM)
      Conversely, cells diverted from terminal differentiation to become long-lived memory B cells possess a quiescent morphology superficially resembling naive B cells. However, they bear the genetic scars of class switching and affinity maturation, expressing specialized, high-affinity isotypes (such as IgG, IgA, or IgE) rather than the baseline IgM/IgD phenotype restricted to naive cohorts.

      % High levels of BSAP (differentiating related factor) tend to promote memory cell formation
      The stabilization and maintenance of the memory cell lineage is highly dependent on the continued, elevated expression of the BSAP transcription factor.

    \end{core}

  \subsection{Leukemia \& Lymphoma}
    
    \begin{core}{Leukemia \& Lymphoma}

      % NON-PRODUCTIVE REARRANGEMENTS
      The Burden of Non-Productive Rearrangements

      % The joining process of the V(D)J regions contains so many random elements, thus
      The highly error-prone and structurally randomized nature of the V(D)J recombination process introduces immense genomic instability.

      % any one junction could produce a frameshift about ⅔ of the time
      Statistically, the imprecise resolution of coding joints possesses roughly a 66\% probability of introducing a catastrophic frameshift mutation during any single recombination event.

      % N and P addition can also produce random stop codons
      Furthermore, the stochastic action of Artemis (P-nucleotides) and TdT (N-nucleotides) frequently generates premature, non-sense stop codons directly within the reading frame.

      % This could all lead to apoptosis because of non-productive rearrangements (you must get both a good H and a good L of some kind)
      Failing to generate structurally viable heavy and light chains results in a non-productive rearrangement, inevitably condemning the developing progenitor cell to immediate apoptosis.

      % ALSO
      % ALLELIC EXCLUSION
      Allelic Exhaustion

      % A cell would rearrange through all H genes (two times) and κ/λ L genes (four times) across the mother and father alleles before it is deemed unproductive and undergoes apoptosis
      A developing cell will exhaust all available allelic options—attempting rearrangement on both heavy chain alleles and across the four total light chain alleles—before finally succumbing to apoptotic clearance.

      % Because of all this gene rearrangement, the developing B cells frequently divide, making blood cancer (leukemia and lymphoma) a higher possibility
      This intense requirement for repeated, double-stranded DNA breaks and rampant cellular proliferation inherently predisposes the B cell lineage to oncogenic transformation. Erroneous recombination events, such as translocations involving the highly active Ig promoters and critical proto-oncogenes, form the fundamental molecular basis for many lymphomas and leukemias.

      <FIGURE>

    \end{core}


